\documentclass[12pt, letterpaper]{article}

% ── Packages ────────────────────────────────────────────────────────────────
\usepackage[margin=1in]{geometry}
\usepackage{amsmath, amssymb}
\usepackage{graphicx}
\usepackage{booktabs}
\usepackage{hyperref}
\usepackage{xcolor}
\usepackage{caption}
\usepackage{subcaption}
\usepackage[authoryear]{natbib}
\usepackage{setspace}
\usepackage{parskip}
\usepackage{microtype}
\microtypesetup{expansion=false}
\usepackage{fancyhdr}
\usepackage{titlesec}
\usepackage{enumitem}
\usepackage{siunitx}
\usepackage{pifont}
\usepackage{newunicodechar}
\newunicodechar{✓}{\checkmark}
\newunicodechar{✗}{\ding{55}}

\hypersetup{
    colorlinks=true,
    linkcolor=blue!60!black,
    citecolor=blue!60!black,
    urlcolor=blue!60!black,
}

\pagestyle{fancy}
\fancyhf{}
\rhead{\small Urban Heat Island Simulation — Austin, TX}
\lhead{\small [Author Name]}
\cfoot{\thepage}

% ── Title ────────────────────────────────────────────────────────────────────
\title{
    \textbf{A Zero-Authentication, Physics-Based Digital Twin for\\
    Urban Heat Island Simulation and Counterfactual\\
    Planning: A Case Study of Austin, Texas}\\[0.5em]
    \large Validated Against MODIS Aqua Land Surface Temperature
}

\author{
    [Author Name]\\
    \small Westwood High School, Round Rock, TX\\
    \small \texttt{[email]}
}

\date{\today}

% ── Document ─────────────────────────────────────────────────────────────────
\begin{document}

\maketitle
\thispagestyle{empty}

% ── Abstract ─────────────────────────────────────────────────────────────────
\begin{abstract}
Urban heat islands (UHIs) pose growing public health risks as cities densify
and climate extremes intensify. Existing high-fidelity simulation tools either
require proprietary meteorological inputs, extensive local calibration, or
significant computational resources, limiting their use in rapid planning
contexts. We present \textsc{austin-twin}, a lightweight, fully open-data
digital twin of Austin, Texas that (1) constructs a 500\,m UTM grid from
OpenStreetMap boundary data and ESA WorldCover 10\,m land-cover tiles accessed
without authentication; (2) simulates 48-hour surface temperature evolution
using a per-cell energy balance with solar absorption, Newtonian radiative
cooling, land-cover-weighted evapotranspiration, mixed-layer water thermal mass,
and horizontal diffusion; (3) calibrates physical coefficients against MODIS
Aqua land surface temperature (LST) observations using differential evolution
with a strict train/test split, using a two-stage procedure that anchors the
horizontal diffusion coefficient to literature-derived tree-cooling decay
scales before refitting absorption and ET coefficients; (4) replaces synthetic
atmospheric forcing with per-date ERA5 hourly reanalysis, enabling absolute
(rather than anomaly-only) surface-temperature comparison within this
framework; and (5) generates counterfactual urban planning scenarios and
spatially targeted sensitivity maps. Relative to existing UHI simulation
systems, the contribution is not a new energy-balance formulation but a
specific combination not previously reported together: a fully
authentication-free multi-sensor remote-sensing pipeline (optical land
cover, thermal-infrared LST, and reanalysis forcing) coupled to a
train/test-disciplined satellite calibration protocol and an explicit,
quantitative diagnostic — a land-cover data-ceiling control — for
separating genuine model failure from sensor-driven class-separability
limits during cross-city transfer. The ship configuration (two-stage
calibration + ERA5 + mixed-layer water physics) achieves a spatial Pearson
correlation of $r = +0.758$ (95\% bootstrap CI, Section~\ref{sec:uncertainty})
and RMSE = \SI{1.42}{\degreeCelsius} on two held-out days, with every
calibrated coefficient in its published physical range and a modeled
absolute spatial spread (\SI{15.4}{\degreeCelsius}) within
\SI{0.4}{\degreeCelsius} of MODIS observations. A diagnostic investigation
of Penman-Monteith ET as a replacement for the linear ET formulation finds
a $\Delta r = -0.33$ regression on the held-out set, attributable in
roughly equal parts to LAI saturation collapsing the vegetation-fraction
spatial gradient and to canonical PM's VPD decoupling eliminating the
surface-temperature feedback that drives spatial UHI contrast; the
mechanism is reported as a documented negative result with implications
for other slab-model ET choices. The framework is entirely reproducible
with a single \texttt{pip install} and no registration requirements,
placing it in an underserved niche between computationally heavy urban
canopy models and correlational machine-learning approaches that lack
causal counterfactual capability. A multi-city transfer experiment on
Phoenix and Miami, interpreted through a satellite land-cover-class
separability lens rather than climate distance, reveals that portability
is determined by whether WorldCover's three-channel representation cleanly
separates a city's thermal regimes: Miami transfers zero-shot at 92\% of
its data ceiling ($r = +0.760$, rising to $r = +0.832$ with three
target-city training days); Phoenix has a low data ceiling
($r = +0.363$) due to irrigation heterogeneity that collapses the
WorldCover-to-temperature mapping, a finding that would be misinterpreted
as model failure without the ceiling control. We report bootstrap
confidence intervals on all reported correlation and RMSE estimates and a
one-at-a-time sensitivity analysis of the calibrated coefficients
(Section~\ref{sec:uncertainty}).
\end{abstract}

\vspace{1em}
\noindent\textbf{Keywords:} land surface temperature retrieval and validation;
MODIS Aqua; multi-sensor data fusion; land-cover class separability; urban
heat island; surface energy balance; digital twin; open remote-sensing data;
counterfactual scenarios; Austin, Texas

\newpage
\tableofcontents
\newpage

% ── 1. Introduction ──────────────────────────────────────────────────────────
\section{Introduction}
\label{sec:intro}

Urban heat islands arise when built-up surfaces — impervious asphalt, concrete,
and rooftops — absorb and retain solar radiation more efficiently than the
vegetation and soil they replace, while simultaneously eliminating the
latent-heat release of transpiration \citep{oke1982energetic}. The resulting
elevated surface and near-surface temperatures carry well-documented
consequences: increased all-cause and heat-related mortality, amplified peak
electricity demand, degraded outdoor air quality, and compounding interaction
with regional drought cycles \citep{voogt2003thermal}.

Austin, Texas presents a particularly demanding test case. Situated in a
semi-arid climate where summer afternoon temperatures routinely exceed
\SI{38}{\degreeCelsius}, the city experienced a sustained heat wave throughout
August 2024, with MODIS Aqua land surface temperatures exceeding
\SI{50}{\degreeCelsius} over downtown impervious surfaces. Rapid suburban
expansion and continued densification of the urban core have widened the
thermal contrast between built and vegetated land covers in recent years.

Two dominant paradigms currently exist for urban heat simulation, and both
carry practical limitations for rapid planning applications. The first
encompasses high-fidelity mesoscale models: the Weather Research and
Forecasting model with urban canopy physics (WRF-SLUCM or WRF-BEP), and
neighborhood-scale tools such as the Urban Weather Generator
\citep{bueno2013urban}, which solve energy conservation across urban canyon
geometries with hourly forcing from operational weather stations. These tools
reproduce observed UHI patterns with high fidelity but require significant
setup: local urban morphology databases, EnergyPlus weather files, and compute
times measured in hours to days for city-scale runs. The second paradigm,
which has seen rapid growth in recent years, applies machine learning —
gradient boosting, CNN-LSTM, and geospatial foundation models — to regress
land surface temperature against spectral indices such as NDVI and NDBI
\citep{lyu2022integrated}. These approaches validate well on historical data
but are inherently correlational: they cannot simulate the temperature effect of
a proposed intervention (``what if we plant trees in this zone?'') because
they contain no causal physical mechanism.

We propose a middle path: a physics-based energy-balance simulator operating
on a 2D spatial grid, driven entirely by open and freely accessible remote-sensing
and reanalysis data sources, with no registration requirements and no local
infrastructure. From a remote-sensing standpoint, the system is best
understood as a sensor-fusion and validation pipeline: it fuses a
thermal-infrared LST product (MODIS Aqua MYD11A1), a 10\,m multispectral
land-cover classification (ESA WorldCover), and coarse-resolution atmospheric
reanalysis (ERA5) into a single per-cell state, calibrates against the LST
product under a strict train/test split, and uses the resulting
satellite-vs-model residuals — rather than model architecture alone — to
diagnose where and why a land-cover-driven thermal model degrades. The
system requires roughly one second of compute per 48-hour simulation on a
consumer laptop, supports city-scale counterfactual scenario generation, and
achieves satellite-validated spatial correlations competitive with approaches
requiring far greater data and setup overhead.

\subsection{Contributions}

\begin{itemize}[noitemsep]
  \item A complete, reproducible pipeline from city boundary to validated
        temperature field using only freely accessible, zero-authentication
        data sources (OSM, ESA WorldCover, MODIS via Microsoft Planetary
        Computer, ERA5 via Copernicus CDS).
  \item A CFL-stable slab energy-balance simulator with configurable
        resolution (500\,m default; 100\,m supported), mixed-layer water
        physics, and a two-stage calibration module using differential
        evolution against multiple MODIS days with a strict train/test split.
  \item A ship configuration (two-stage + ERA5 + mixed-layer water) achieving
        $r = +0.758$, RMSE = \SI{1.42}{\degreeCelsius} on held-out days with
        every calibrated coefficient in its published physical range and
        absolute thermal spread within \SI{0.4}{\degreeCelsius} of MODIS.
  \item A documented negative result for canonical Penman-Monteith ET,
        with the two contributing mechanisms isolated: LAI saturation
        collapsing the vegetation-fraction gradient ($\sim$1/3 of the
        regression) and VPD decoupling eliminating the surface-temperature
        feedback that drives spatial UHI contrast ($\sim$2/3).
  \item A multi-city transfer experiment on Phoenix and Miami with a
        data-ceiling control showing that portability is determined by
        WorldCover land-cover separability rather than climate distance,
        and that Phoenix's low transfer score is a ceiling effect rather
        than model failure.
  \item A counterfactual scenario engine with five canonical urban planning
        interventions and a zonal sensitivity analysis identifying where
        canopy investment yields maximum cooling per unit area.
\end{itemize}

\subsection{What Is and Is Not New}
\label{sec:novelty}

None of the individual pieces of this system is, in isolation, a new idea:
slab-style urban energy balances predate this work \citep{oke1982energetic},
differential-evolution calibration against satellite LST is standard
practice \citep{mitraka2025prototypes}, and district-scale material-driven
thermal simulators validated against optical/thermal satellite imagery
already exist \citep{heatmat2025}. We do not claim a new energy-balance
formulation, a new calibration algorithm, or a new satellite product.
Table~\ref{tab:novelty} makes explicit what combination of design choices we
believe is new relative to the four closest reference points — mesoscale
urban canopy models (WRF-SLUCM/BEP, TEB), the Urban Weather Generator,
correlational ML approaches, and the most closely related digital-twin
system, HeatMat.

\begin{table}[ht]
  \centering
  \caption{Positioning relative to existing UHI simulation and prediction
  systems. ``Zero-auth data'' means every input can be retrieved by an
  anonymous HTTP request with no API key or account. ``Counterfactual''
  means the system can evaluate an intervention outside the training/forcing
  distribution via an explicit causal mechanism, not interpolation.}
  \label{tab:novelty}
  \small
  \begin{tabular}{lccccc}
    \toprule
    System & Spatial scale & Zero-auth data & Counterfactual & Sat. calib. & Ceiling diagnostic \\
    \midrule
    WRF-SLUCM / TEB     & Meso / canyon   & \ding{55} & \checkmark & rare        & \ding{55} \\
    Urban Weather Gen.\ & Neighborhood    & \ding{55} & \checkmark & \ding{55}   & \ding{55} \\
    ML (NGBoost/CNN-LSTM/GeoFM) & City--regional & partial & \ding{55} & implicit & \ding{55} \\
    HeatMat (2025)      & District        & partial & \ding{55}  & \checkmark  & \ding{55} \\
    \textbf{This work}  & City            & \checkmark & \checkmark & \checkmark  & \checkmark \\
    \bottomrule
  \end{tabular}
\end{table}

Two entries in Table~\ref{tab:novelty} carry the actual novelty claim of this
paper. First, \emph{zero-auth data at city scale combined with a
counterfactual engine}: WRF-class models and the Urban Weather Generator
support counterfactuals but require registered meteorological products or
manually digitized morphology; ML approaches use openly accessible imagery
but cannot evaluate interventions outside their training distribution; to
our knowledge no existing system combines fully anonymous-access data
retrieval, city-scale (not district-scale) coverage, and an explicit physical
counterfactual mechanism. Second, and more central to this venue, the
\emph{ceiling diagnostic}: none of the four reference systems reports a
land-cover-class-separability upper bound alongside transfer results, which
means a low cross-city correlation in those frameworks is not distinguishable
from model failure. Section~\ref{sec:transfer} shows this distinction is not
cosmetic — it reverses the interpretation of the Phoenix result. We consider
this diagnostic, not the energy-balance model itself, the paper's most
transferable methodological contribution.

\subsection{Relevance to Remote Sensing}
\label{sec:rs_relevance}

The primary evidentiary content of this paper is a set of satellite-vs-model
comparisons: every quantitative claim in Sections~\ref{sec:results}
and~\ref{sec:transfer} is a statement about how well a physical model,
forced and validated by remote-sensing products, reproduces an independently
sensed thermal-infrared field. The methodological questions we address —
how to fuse optical land-cover classification, thermal-infrared LST, and
reanalysis forcing at mismatched native resolutions ($10\,\text{m}$,
$1\,\text{km}$, and $\sim$$30\,\text{km}$ respectively); how to attribute a
degraded satellite-model correlation to sensor/classification limitations
versus model error; and how QA-flag masking and per-date overpass timing
propagate into calibrated physical parameters — are remote-sensing data
fusion and validation questions rather than urban-planning questions. The
counterfactual scenario engine (Section~\ref{sec:results}) is a downstream
application of the validated sensor-model fusion, not the paper's central
technical contribution.

% ── 2. Related Work ──────────────────────────────────────────────────────────
\section{Related Work}
\label{sec:related}

\subsection{Physics-Based Urban Energy Balance Models}

The Town Energy Balance (TEB) scheme \citep{masson2000physically} and
the Single-Layer Urban Canopy Model \citep{kusaka2001simple} embedded in
WRF represent the physics-heavy end of the spectrum, resolving solar geometry,
longwave radiation, and surface conductance across explicit street-canyon
geometries. The Urban Weather Generator \citep{bueno2013urban} provides a
lighter alternative that estimates urban canopy air temperature from rural
weather-station data via neighborhood-scale energy balances, achieving
satisfactory performance in Basel, Toulouse, and Abu Dhabi while remaining
fast enough for iterative design workflows. Our slab model is lighter still:
it sacrifices canyon geometry and 3D radiation transfer in exchange for
city-scale 2D spatial fields, computational speed, and zero dependence on
local meteorological infrastructure.

A recent comparable effort, \textit{HeatMat} \citep{heatmat2025}, simulates
urban material-driven surface temperatures at district scale and validates
against Landsat imagery in Detroit. Our work differs primarily in scope
(city-scale vs.\ district-scale), data pipeline (fully automated from
OSM/WorldCover with no manual digitization), and the inclusion of a
counterfactual intervention engine.

\subsection{Machine Learning Approaches}

Recent work has increasingly applied machine learning to UHI prediction.
\citet{lyu2022integrated} combine satellite imagery with urban sensor networks
in an integrated CyberGIS-ML framework for fine-scale UHI prediction.
\citet{hoang2025geospatial} apply natural gradient boosting and deep neural
networks to multisource remote sensing data in the Midwest, achieving strong
cross-validated correlations. Foundation models pre-trained on geospatial data
have also been applied to LST detection \citep{bhamjee2024detection}.
These approaches can achieve high predictive accuracy on held-out data,
but because they learn statistical associations between land-cover indices
and observed temperatures, they cannot be used to evaluate the effect of a
proposed intervention that lies outside the training distribution —
a central use case for urban planning.

\subsection{Satellite Validation of UHI Models}

MODIS Aqua LST (product MYD11A1, \SI{1}{\km} nominal resolution) has been
used extensively as a validation target for urban heat models due to its
long archive, frequent overpass, and publicly accessible retrieval via
platforms including Microsoft Planetary Computer and NASA Earthdata. The
standard practice of anomaly-based correlation — subtracting the citywide
spatial mean from both model and observation before computing Pearson $r$ —
allows validation of spatial pattern even when absolute temperature offsets
are present due to forcing differences \citep{mitraka2025prototypes}.
We adopt this approach and additionally report absolute RMSE after calibration
closes most of the mean offset.

% ── 3. Data ──────────────────────────────────────────────────────────────────
\section{Data}
\label{sec:data}

All data sources used in this study are publicly accessible without
registration or authentication, supporting full reproducibility.

\subsection{City Boundary}
The Austin city-limits polygon was obtained from OpenStreetMap via the
\texttt{osmnx.geocode\_to\_gdf} function and cached locally as GeoJSON.
The polygon defines all spatial masking operations in the pipeline.

\subsection{Land Cover}
ESA WorldCover v200 (2021 reference year) provides 10\,m global land cover
in 11 classes \citep{zanaga2022esa}. Tiles covering Austin were accessed
as cloud-optimized GeoTIFFs (COGs) from the public AWS S3 bucket
\texttt{esa-worldcover}, requiring no authentication. Tiles are cached
locally on first download and reused across all subsequent runs.

Three continuous channels were derived by aggregating WorldCover pixels
within each simulator grid cell:
\begin{itemize}[noitemsep]
  \item \textbf{impervious\_frac}: fraction of pixels classified as
        Built-up (class 50)
  \item \textbf{vegetation\_frac}: fraction classified as Tree, Shrub,
        Grass, Crop, Wetland, or half of Bare (classes 10, 20, 30, 40, 90,
        plus 0.5$\times$60)
  \item \textbf{water\_mask}: fraction classified as Permanent Water
        (class 80)
\end{itemize}

\subsection{Surface Temperature}
MODIS Aqua MYD11A1 v6.1 daily land surface temperature was retrieved via
the Microsoft Planetary Computer STAC API, which provides open access to
the full archive with no authentication \citep{microsoft_pc}. The Aqua
overpass time over Austin is approximately 1:30\,pm local time, corresponding
to peak or near-peak daytime surface temperatures. Pixels with QA flags
indicating cloud contamination or retrieval errors were masked.

For Austin, five overpass days in August 2024 were used: August 11, 16, and 19
(calibration/training) and August 21 and 22 (held-out test). All days fell
within the 2024 Texas heat wave.

For the multi-city transfer experiment (Section~\ref{sec:transfer}), MODIS
days were selected to provide warm-season coverage representative of each
city's peak UHI regime: July 2024 for Phoenix (peak desert heat wave) and
April 2024 for Miami (dry season, high surface insolation). Three days
per city were used for any adaptation calibration; two held-out days per city
were used exclusively for evaluation. No held-out days from any target city
entered any calibration step.

\subsection{Atmospheric Forcing}
ERA5 hourly reanalysis data were retrieved from the Copernicus Climate Data
Store for each city and date window: August 10--23, 2024 for Austin;
July 15--28, 2024 for Phoenix; and April 4--20, 2024 for Miami. Variables
retrieved: 2\,m air temperature (\texttt{t2m}), 2\,m dewpoint temperature
(\texttt{d2m}), surface solar radiation downwards (\texttt{ssrd}), and 10\,m
wind components (\texttt{u10}, \texttt{v10}). ERA5 is produced by the
European Centre for Medium-Range Weather Forecasts (ECMWF) and is freely
available under a CC-BY 4.0 license \citep{hersbach2020era5}. \texttt{ssrd}
is an accumulated variable and was de-accumulated by differencing consecutive
hourly values before use. Each city's ERA5 data were cached to a separate
named NetCDF to prevent cross-city data corruption (a bug in an earlier version
that used a shared unpack directory; see Section~\ref{sec:transfer}).

% ── 4. Methods ───────────────────────────────────────────────────────────────
\section{Methods}
\label{sec:methods}

\subsection{Spatial Grid}

The simulation domain is a regular grid in UTM Zone 14N (EPSG:32614),
clipped to the Austin city-limits polygon. The default cell size is 500\,m,
yielding an approximately $94 \times 73$ array with $\sim$2,896 valid city
cells. A higher-resolution 100\,m grid is also supported via an environment
variable (\texttt{AUSTIN\_GRID\_M=100}), producing $\sim$72,486 valid cells
and resolving highway-corridor-scale features at the cost of approximately
25$\times$ longer simulation time.

\subsection{Slab Energy Balance Simulator}

Each grid cell evolves independently under a per-cell surface energy balance:
\begin{equation}
  \frac{dT_i}{dt} = Q_{\mathrm{solar}}(t, \mathbf{L}_i)
                  - Q_{\mathrm{rad}}(T_i)
                  - Q_{\mathrm{ET}}(T_i, \mathbf{L}_i)
                  + Q_{\mathrm{diff}}(T, i)
  \label{eq:energy_balance}
\end{equation}
where $T_i$ is the surface temperature of cell $i$, $\mathbf{L}_i$ denotes
its land-cover composition vector (impervious, vegetation, and water fractions),
and $t$ is simulation time.

\paragraph{Solar absorption.}
Incoming solar radiation follows a half-sine diurnal profile with peak at
solar noon and zero at night:
\begin{equation}
  Q_{\mathrm{solar}}(t, \mathbf{L}_i) = I_{\max}(t) \cdot
  \bigl(\alpha_{\mathrm{imp}} \cdot f_{\mathrm{imp},i}
       + \alpha_{\mathrm{veg}} \cdot f_{\mathrm{veg},i}
       + \alpha_{\mathrm{wat}} \cdot f_{\mathrm{wat},i}\bigr)
\end{equation}
where $\alpha_{\mathrm{imp}}$, $\alpha_{\mathrm{veg}}$, and
$\alpha_{\mathrm{wat}}$ are the per-class effective solar absorption
coefficients (calibrated parameters) and $I_{\max}(t)$ is the diurnal solar
forcing amplitude.

\paragraph{Radiative cooling.}
Longwave radiative loss is approximated as Newtonian cooling toward a
background ambient temperature $T_{\mathrm{air}}$:
\begin{equation}
  Q_{\mathrm{rad}}(T_i) = k_{\mathrm{rad}} \cdot (T_i - T_{\mathrm{air}})
\end{equation}

\paragraph{Evapotranspiration.}
Latent heat removal by evapotranspiration is modeled as proportional to
vegetation fraction and the excess of surface temperature above ambient,
one-sided (no condensation):
\begin{equation}
  Q_{\mathrm{ET}}(T_i, \mathbf{L}_i) =
  k_{\mathrm{ET}} \cdot f_{\mathrm{veg},i} \cdot
  \max(T_i - T_{\mathrm{air}}, 0)
\end{equation}

\paragraph{Horizontal diffusion.}
Thermal coupling between adjacent cells is modeled by a five-point Laplacian
operator with reflective boundary conditions:
\begin{equation}
  Q_{\mathrm{diff}}(T, i) = D \cdot \nabla^2 T \big|_i
\end{equation}
The timestep $\Delta t$ is set automatically to satisfy the CFL stability
condition $\Delta t \leq \Delta x^2 / (4D)$ at all resolutions.

\paragraph{Water cells.}
\label{par:water_methods}
Cells with substantial water fraction relax toward the diurnal mean ambient
temperature on a timescale proportional to $f_{\mathrm{wat},i}$, approximating
the high thermal mass of surface water without a full deep-layer model.

\paragraph{Atmospheric forcing.}
\label{par:forcing}
The ship configuration uses per-date ERA5 hourly reanalysis fields
(Section~\ref{sec:data}) rather than a synthetic diurnal profile. For each
simulation date, the ERA5 nearest-grid-point time series for Austin
(\texttt{t2m}, \texttt{ssrd}) is interpolated to the simulator timestep and
used directly as $T_{\mathrm{air}}(t)$ and $I_{\max}(t)$. This eliminates
the synthetic-forcing assumption that previously required all validation
within this project to be performed in anomaly space (subtracting the
citywide spatial mean before correlation), and allows a direct comparison
of modeled and observed absolute surface temperatures rather than only
spatial anomalies. When ERA5 forcing is active, the simulator initializes
each grid cell at the temporal mean of the forcing $T_{\mathrm{air}}$
time series rather than the fixed synthetic-forcing default of
30\,\si{\degreeCelsius}; the latter was found to cause a systematic
spin-up artifact on cold-season days when the default deviates
substantially from the actual ambient temperature, manifesting as an
incorrect spatial pattern during the first 24--48\,h of simulation rather
than the forcing-driven pattern. A synthetic 30\,$\pm$\,6\,\si{\degreeCelsius}
profile is retained as the forcing for the single-stage baseline configuration
and for the canopy-sweep and zonal-sensitivity analyses, which must be
run under fixed forcing to ensure counterfactual comparisons are attributed
only to the land-cover change.

\subsection{Calibration}
\label{sec:calibration}

We implement and compare two calibration strategies that differ in what
physical objective $D$ is asked to satisfy.

\paragraph{Stage 1 (single-stage): baseline MODIS pattern.}
Five simulator coefficients are jointly optimized against MODIS observations:
\begin{equation}
  \boldsymbol{\theta} = \{
    \alpha_{\mathrm{imp}},\;
    \alpha_{\mathrm{veg}},\;
    \alpha_{\mathrm{wat}},\;
    k_{\mathrm{ET}},\;
    D
  \}
\end{equation}
The objective minimizes mean anomaly RMSE across three training MODIS days
(August 11, 16, and 19, 2024):
\begin{equation}
  \hat{\boldsymbol{\theta}} = \arg\min_{\boldsymbol{\theta}}
  \frac{1}{N_{\mathrm{train}}} \sum_{d \in \mathcal{D}_{\mathrm{train}}}
  \mathrm{RMSE}\bigl(
    \tilde{T}^{\mathrm{sim}}_d(\boldsymbol{\theta}),\;
    \tilde{T}^{\mathrm{MODIS}}_d
  \bigr)
\end{equation}
where $\tilde{T}$ denotes the spatially de-meaned anomaly field.
Optimization uses \texttt{scipy.optimize.differential\_evolution}
\citep{storn1997differential}. Two held-out days (August 21 and 22) are
scored only after optimization completes and never enter the loss function.
This strategy produces the best achievable MODIS pattern score but allows $D$
to take whatever value minimizes RMSE — which, as discussed in
Section~\ref{par:diffusion}, may not correspond to physically realistic
tree-cooling decay scales.

\paragraph{Stage 2 (two-stage): decay-constrained diffusion.}
To disentangle baseline coherence from intervention spread, a second
calibration anchors $D$ to the published canopy-cooling literature before
re-fitting the absorption and ET coefficients. Specifically: (i) a sweep
of candidate $D$ values is run with all other coefficients fixed at their
single-stage values; (ii) the value of $D$ that minimizes the discrepancy
between the simulated patch half-cooling distance and the midpoint of the
empirical literature band ($\sim$$200\,\text{m}$) is selected; (iii) with
$D$ now fixed, the four absorption and ET coefficients
$\{\alpha_{\mathrm{imp}}, \alpha_{\mathrm{veg}}, \alpha_{\mathrm{wat}},
k_{\mathrm{ET}}\}$ are re-optimized against the same three MODIS training
days. Both configurations are retained: single-stage for applications
where baseline MODIS fidelity is the primary criterion, two-stage for
applications where physically realistic intervention-effect spatial
distributions are required.

\paragraph{Ship configuration: two-stage + ERA5.}
\label{par:ship}
The \emph{ship configuration} used for primary quantitative results combines
the two-stage diffusion anchor ($D = 1.0\,\text{m}^2\,\text{s}^{-1}$,
pinned) with per-date ERA5 atmospheric forcing (Section~\ref{par:forcing})
and mixed-layer water physics (Section~\ref{par:water_methods}).
Calibration re-fits $\{\alpha_{\mathrm{imp}}, \alpha_{\mathrm{veg}},
\alpha_{\mathrm{wat}}, k_{\mathrm{ET}}\}$ against the same three training
MODIS days using ERA5 forcing for each date, with $D$ pinned at
$1.0\,\text{m}^2\,\text{s}^{-1}$. The decision not to adopt a
Penman-Monteith ET formulation in the ship configuration is explained by
a diagnostic experiment described in Section~\ref{sec:pm_negative}.

\subsection{Counterfactual Scenarios}

Each scenario is implemented as a callable transformation of the land-cover
dataset. Five canonical scenarios are provided:
\begin{enumerate}[noitemsep]
  \item \textbf{Baseline}: unmodified WorldCover land cover.
  \item \textbf{Canopy +20\%}: vegetation fraction increased by 0.20 for all
        cells, impervious fraction reduced correspondingly.
  \item \textbf{Cool roofs downtown}: impervious absorption coefficient
        reduced by 0.35 within a 3\,km radius of the city center, representing
        high-albedo roof coatings.
  \item \textbf{River greenway}: a 1\,km corridor along Barton Creek and
        the Colorado River converted from impervious to vegetation.
  \item \textbf{Suburban densification}: outer-ring residential cells
        converted to higher impervious fraction, reflecting continued growth.
\end{enumerate}

\subsection{Zonal Sensitivity Analysis}

To identify where canopy investment yields the most cooling per unit area,
the city is partitioned into a $6 \times 5$ zone grid ($\sim$27 zones with
$\geq 6$ valid cells after boundary trimming). For each zone, a local canopy
scenario (vegetation fraction $+0.20$ inside the zone only) is simulated,
and we compute: local mean $\Delta T$, citywide mean $\Delta T$, total
cooling ($\si{\degreeCelsius\cdot m^2}$), and cooling efficiency
(total cooling per unit converted area). Zones with the highest efficiency
represent the best candidates for targeted investment.

\subsection{Multi-City Transfer Experiment}
\label{sec:transfer_methods}

To assess geographic generalizability, the ship configuration is evaluated
on two target cities — Phoenix, AZ and Miami, FL — chosen to span orthogonal
climate and land-cover axes. Phoenix is climatically similar to Austin
(semi-arid, hot summers, high solar irradiance) but has a distinct urban
morphology dominated by low-density sprawl, extensive irrigated landscaping,
and a high fraction of bare desert soil within WorldCover's ``vegetation''
class. Miami is climatically distant (subtropical coastal, high humidity,
tropical vegetation) but has sharply differentiated WorldCover land-cover
categories with strong thermal contrast between impervious surfaces and
genuine dense tropical canopy.

Three transfer schemes are evaluated for each target city:

\begin{itemize}[noitemsep]
  \item \textbf{Zero-shot}: Austin ship-configuration coefficients applied
        directly, with no target-city data.
  \item \textbf{Scheme A (full re-calibration)}: All four coefficients
        $\{\alpha_{\mathrm{imp}}, \alpha_{\mathrm{veg}}, \alpha_{\mathrm{wat}},
        k_{\mathrm{ET}}\}$ re-optimized against 1 or 3 target-city MODIS days
        via differential evolution.
  \item \textbf{Scheme B (absorption freeze)}: Only $k_{\mathrm{ET}}$ is
        re-optimized against 1 or 3 target-city MODIS days, treating absorption
        coefficients as physics-portable and only adjusting the ET coupling
        to local climate.
\end{itemize}

For each scheme and city, two held-out MODIS days (never seen during any
calibration step) are used to compute test-set Pearson $r$.

To interpret the results, we compute a \textbf{data-ceiling control}:
the Pearson correlation between MODIS LST and impervious fraction across
all valid city cells. This is the maximum correlation any land-cover-only
model could achieve regardless of calibration, and provides a physically
motivated upper bound against which transfer performance is assessed.
A model scoring at or near the ceiling has extracted all available
land-cover information; a model scoring well below it has failed to do so.

\subsection{Uncertainty Quantification and Sensitivity Analysis}
\label{sec:uncertainty_methods}

Point estimates of $r$ and RMSE reported elsewhere in this paper are
computed from two held-out MODIS days per city; without an estimate of
sampling and spatial-autocorrelation uncertainty, differences of a few
hundredths in $r$ (e.g.\ the $\Delta r = +0.009$ attributed to ERA5 forcing
in Section~\ref{sec:ship_results}) cannot be distinguished from noise. We
therefore report interval estimates rather than point estimates alone, and
a sensitivity analysis of the calibrated coefficients.

\paragraph{Spatial block bootstrap for $r$ and RMSE.}
Because the standard closed-form CI for a Pearson correlation (Fisher's
$z$-transform) assumes independent observations, and adjacent grid cells
in this model are spatially correlated by construction (both physically,
via the diffusion term, and observationally, via MODIS's $\sim$1\,km
footprint and the model's own $500$\,m grid), we use a \emph{spatial block
bootstrap} rather than a naive per-cell bootstrap, which would understate
the true interval width. Procedure: (i) partition the valid-cell grid into
contiguous $L \times L$ blocks ($L = 6$ cells, $\approx 3$\,km, chosen to
exceed the fitted horizontal correlation length of the diffusion field);
(ii) resample blocks with replacement until the resampled set has
approximately the same number of cells as the original grid; (iii) recompute
Pearson $r$ and RMSE between simulator and MODIS on the resampled field;
(iv) repeat for $B = 2{,}000$ bootstrap replicates per test day; (v) pool
replicates across the held-out test days for that configuration and report
the 2.5th/97.5th percentiles as the 95\% CI. This is applied independently
to every $(r, \mathrm{RMSE})$ pair reported for a held-out test set in
Sections~\ref{sec:cal_results}, \ref{sec:ship_results}, and
\ref{sec:transfer}.

\paragraph{Coefficient sensitivity.}
For the ship configuration's four re-fitted coefficients
$\{\alpha_{\mathrm{imp}}, \alpha_{\mathrm{veg}}, \alpha_{\mathrm{wat}},
k_{\mathrm{ET}}\}$ (with $D$ held pinned, per Section~\ref{par:ship}), we
report a local one-at-a-time (OAT) sensitivity: each coefficient is
perturbed by $\pm 10\%$ and $\pm 25\%$ from its calibrated value with all
others held fixed, and the resulting held-out test $r$ and RMSE are
recomputed. This is a local, not global, sensitivity measure — it does not
capture coefficient interactions — but is sufficient to establish which
coefficients the ship configuration's held-out skill is most fragile to,
and to give the calibrated values an approximate confidence radius rather
than treating them as exact. We additionally report the spread of the
final differential-evolution population (not just its best member) as a
complementary, optimization-native measure of how tightly each coefficient
is constrained by the training-day objective.

\paragraph{Scope and honest limitations of this analysis.}
With only two held-out days per city, block-bootstrap CIs describe
within-day spatial sampling uncertainty; they do not capture
day-to-day (synoptic) variability, which Section~\ref{sec:extended_validation}'s
18-day extended validation addresses separately and which is the more
relevant uncertainty source for claims about generalization across dates.
The two sources are complementary: the bootstrap CI answers ``how
precisely do we know $r$ on \emph{this} day's spatial field,'' while the
extended-validation spread answers ``how much does $r$ vary from
\emph{day to day}.'' We report both and do not conflate them.

% ── 5. Results ───────────────────────────────────────────────────────────────
\section{Results}
\label{sec:results}

\subsection{Uncalibrated Baseline Validation}

Prior to any calibration, the slab model was run with hand-selected
parameters. Comparison against MODIS Aqua LST on August 19, 2024 yielded
a spatial Pearson correlation of $r = +0.788$ with an anomaly RMSE of
\SI{1.71}{\degreeCelsius} (Figure~\ref{fig:uncalibrated}). This result
demonstrates that the pattern correlation — the fundamental spatial
structure of the UHI, with impervious hotspots and vegetated cool zones
in the correct locations — emerged directly from the physics and land-cover
data, without any fitting to satellite observations.

The simulator anomaly spread at this stage was approximately
\SI{8.2}{\degreeCelsius} peak-to-peak, compared to approximately
\SI{12}--\SI{16}{\degreeCelsius} in the MODIS observations, indicating
that the physics correctly identifies which cells are hot and which are
cool, but underestimates the magnitude of contrast.

\begin{figure}[ht]
  \centering
  % Replace with actual image path:
  % \includegraphics[width=\linewidth]{figures/validation_uncalibrated.png}
  \fbox{\parbox{0.9\linewidth}{\centering\vspace{4cm}
    [Figure: Uncalibrated validation on Aug 19 2024.\\
    Left: MODIS Aqua LST; Right: Simulator T (synthetic forcing).\\
    Bottom row: anomaly fields. $r = +0.788$, RMSE = 1.71 °C.]\vspace{4cm}}}
  \caption{Uncalibrated simulator vs.\ MODIS Aqua LST on August 19, 2024.
  Both fields are shown in absolute temperature (top) and as anomalies
  from the citywide spatial mean (bottom). The simulator uses synthetic
  diurnal forcing, so absolute values differ; comparison is in anomaly
  space. Pearson $r = +0.788$, anomaly RMSE = \SI{1.71}{\degreeCelsius}
  over 2,896 jointly valid cells.}
  \label{fig:uncalibrated}
\end{figure}

\subsection{Calibration Results}
\label{sec:cal_results}

\paragraph{Single-stage.}
Differential evolution converged on coefficients that improve mean anomaly
RMSE on training days from \SI{1.465}{\degreeCelsius} to
\SI{1.256}{\degreeCelsius}, with mean Pearson $r$ rising from $+0.784$ to
$+0.823$. The held-out test set improved \emph{more} than the training set
(Table~\ref{tab:calibration}), confirming the calibration discovered
physically meaningful structure rather than overfitting to training noise.

\begin{table}[ht]
  \centering
  \caption{Single-stage calibration performance. ``Before'' = hand-selected
  initial coefficients; ``After'' = differential evolution result, evaluated
  under the original legacy water-buffer parameterization.
  (When re-evaluated under current mixed-layer water physics, the test
  metrics shift marginally to $r = +0.850$, RMSE = \SI{1.09}{\degreeCelsius};
  see Table~\ref{tab:two_stage}.)}
  \label{tab:calibration}
  \begin{tabular}{lccccc}
    \toprule
    Split & Days & RMSE before & RMSE after & $r$ before & $r$ after \\
    \midrule
    Train & Aug 11, 16, 19
      & \SI{1.465}{\degreeCelsius} & \SI{1.256}{\degreeCelsius}
      & $+0.784$ & $+0.823$ \\
    Test (held-out) & Aug 21, 22
      & \SI{1.304}{\degreeCelsius} & \SI{1.070}{\degreeCelsius}
      & $+0.814$ & $+0.855$ \\
    \bottomrule
  \end{tabular}
\end{table}

Calibration expanded the simulator's anomaly spread from
\SI{8.2}{\degreeCelsius} to \SI{11.4}{\degreeCelsius}, closing approximately
half the remaining magnitude gap against MODIS. The calibrated result on
August 19 ($r = +0.829$, RMSE = \SI{1.41}{\degreeCelsius}) is shown in
Figure~\ref{fig:calibrated}.

\begin{figure}[ht]
  \centering
  % \includegraphics[width=\linewidth]{figures/validation_calibrated.png}
  \fbox{\parbox{0.9\linewidth}{\centering\vspace{3.5cm}
    [Figure: Calibrated validation on Aug 19 2024.\\
    Left: MODIS anomaly; Center: calibrated simulator anomaly;\\
    Right: per-cell scatter ($r = +0.829$, RMSE = 1.41 °C).]\vspace{3.5cm}}}
  \caption{Single-stage calibrated simulator vs.\ MODIS Aqua LST on August 19,
  2024 (a training day). Left: MODIS anomaly. Center: simulator anomaly.
  Right: per-cell scatter with identity line.
  $r = +0.829$, RMSE = \SI{1.41}{\degreeCelsius}. The held-out test days
  (Aug 21–22) achieve mean $r = +0.850$, RMSE = \SI{1.09}{\degreeCelsius}.}
  \label{fig:calibrated}
\end{figure}

\paragraph{Two-stage.}
Anchoring $D$ to a target half-cooling distance of $\sim$200\,m
(literature midpoint) and re-optimizing the four absorption and ET
coefficients produces a substantially different coefficient set
(Table~\ref{tab:coef_comparison}) and a quantified trade-off between
baseline MODIS pattern fidelity and intervention spatial realism
(Table~\ref{tab:two_stage}).

\begin{table}[ht]
  \centering
  \caption{Calibrated coefficients: single-stage vs.\ two-stage, with
  literature reference ranges where available.}
  \label{tab:coef_comparison}
  \begin{tabular}{lccc}
    \toprule
    Coefficient & Single-stage & Two-stage & Literature ref. \\
    \midrule
    $\alpha_{\mathrm{imp}}$ & 0.827 & 0.800 & $\sim$0.85--0.95 \\
    $\alpha_{\mathrm{veg}}$ & 0.371 & 0.498 & $\sim$0.50--0.65 \\
    $\alpha_{\mathrm{wat}}$ & 0.119 & 0.169 & $\sim$0.06--0.10 \\
    $k_{\mathrm{ET}}$ (W\,m$^{-2}$\,K$^{-1}$\,veg$^{-1}$) & 13.25 & 6.85 & broad \\
    $D$ (m$^2$\,s$^{-1}$) & 18.31 & 1.00 & anchored to decay \\
    \midrule
    Patch half-cooling distance & 619\,m & 363\,m & 100--300\,m \\
    \bottomrule
  \end{tabular}
\end{table}

\begin{table}[ht]
  \centering
  \caption{MODIS pattern fidelity vs.\ intervention spatial realism
  trade-off between the two calibration strategies.}
  \label{tab:two_stage}
  \begin{tabular}{lcc}
    \toprule
    Metric & Single-stage & Two-stage \\
    \midrule
    Test $r$ (Aug 21--22, held-out)   & $+0.850$ & $+0.747$ \\
    Test RMSE (°C)                     & 1.09     & 1.48     \\
    Patch half-cooling distance        & 619\,m   & 363\,m   \\
    $\alpha_{\mathrm{veg}}$ vs.\ literature & $2\times$ low & in range \\
    $\alpha_{\mathrm{wat}}$ vs.\ literature & in range & in range ✓ \\
    \bottomrule
  \end{tabular}
\end{table}

The two-stage result has two clear gains and one residual problem.
On the gain side: $\alpha_{\mathrm{veg}}$ snaps from an unphysically
low 0.37 to 0.50, squarely within the published range for mixed urban
vegetation (0.50--0.65), indicating the single-stage value was a
compensatory artifact of over-diffusion; $\alpha_{\mathrm{wat}}$ returns
to a physically realistic 0.17 (literature: 0.06--0.10, within a factor
of two), correcting the red-flag value of 0.49 that appeared in a
preliminary run before mixed-layer water physics was implemented; and the
patch half-cooling distance halves from 619\,m to 363\,m, reducing the
gap against the literature band (100--300\,m) by roughly 60\%.
On the residual problem side: the patch half-cooling did not fully reach
the target band. With $D$ anchored using the single-stage absorptions,
the subsequent refitting of absorptions to sharper material contrasts
extended the visible cooling slightly beyond the Stage-1 prediction of
247\,m. A second iteration (refit $D$ given the new absorptions) would
close this gap but is not pursued here.

The two-stage result still outperforms the land-cover-only regression
baseline on held-out MODIS days ($r = +0.747$ vs.\ $+0.692$) and
approaches near-parity with the position-augmented regression at
$+0.763$, so the physics advantage, though narrowed, persists under the
more constrained coefficient set.

\paragraph{Ship configuration: two-stage + ERA5.}
\label{sec:ship_results}
Adding per-date ERA5 atmospheric forcing to the two-stage configuration
(the ship configuration; Section~\ref{par:ship}) yields calibrated
coefficients $\alpha_{\mathrm{imp}} = 0.803$, $\alpha_{\mathrm{veg}} =
0.434$, $\alpha_{\mathrm{wat}} = 0.178$, $k_{\mathrm{ET}} = 10.975$,
$D = 1.000$ (pinned). Per-day and aggregate metrics are given in
Table~\ref{tab:ship_metrics}.

\begin{table}[ht]
  \centering
  \caption{Ship configuration (two-stage + ERA5 + mixed-layer water)
  calibration and held-out test metrics. All scored in absolute
  temperature space — no anomaly normalization required.
  Source: \texttt{outputs/calibrated\_config\_ship.json}.}
  \label{tab:ship_metrics}
  \begin{tabular}{lcccccc}
    \toprule
    Split & Date & $r$ before & $r$ after & RMSE before & RMSE after
           & Sim spread (°C) \\
    \midrule
    Train & 2024-08-11 & $+0.693$ & $+0.713$ & 1.47 & 1.49 & 14.30 \\
    Train & 2024-08-16 & $+0.704$ & $+0.723$ & 1.55 & 1.54 & 14.78 \\
    Train & 2024-08-19 & $+0.743$ & $+0.761$ & 1.65 & 1.61 & 15.39 \\
    \midrule
    Train mean & & $+0.714$ & $+0.732$ & 1.56 & 1.55 & \\
    \midrule
    Test & 2024-08-21 & $+0.719$ & $+0.744$ & 1.52 & 1.51 & 15.34 \\
    Test & 2024-08-22 & $+0.750$ & $+0.772$ & 1.36 & 1.34 & 14.86 \\
    \midrule
    \textbf{Test mean} & & $+0.735$ & $\mathbf{+0.758}$
                       & 1.44 & $\mathbf{1.42}$ & \\
    \bottomrule
  \end{tabular}
\end{table}

The ship configuration achieves test $r = +0.758$ and
RMSE = \SI{1.42}{\degreeCelsius}, a marginally positive improvement of
$\Delta r = +0.011$ over the two-stage synthetic baseline ($r = +0.747$) at
essentially zero cost in RMSE. Although the marginal bootstrap CIs for the
two configurations overlap heavily (Table~\ref{tab:bootstrap_ci}), a
paired-difference bootstrap on the same resampled block indices — which
isolates the ERA5 contribution net of shared spatial-sampling noise —
yields paired $\Delta r = +0.0105$, 95\% CI [$+0.0051$, $+0.0160$]:
the ERA5 gain is statistically distinguishable from zero at the 95\% level
(Section~\ref{sec:uncertainty}). We describe it as marginally positive but
small in absolute terms; its principal significance is methodological rather
than a substantial correlation gain. Critically, the
ship configuration operates in absolute temperature space: on August 19,
2024, the simulator produces a mean surface temperature of
$\SI{40.77}{\degreeCelsius}$ against a MODIS mean of
$\SI{47.56}{\degreeCelsius}$ — a mean offset of $\SI{6.8}{\degreeCelsius}$
attributable to the model's simplified radiative scheme rather than forcing
— with a spatial spread of $\SI{15.39}{\degreeCelsius}$ against MODIS
$\SI{15.82}{\degreeCelsius}$, close to the observed city-scale thermal
range. Within this project's development history, this is the first
configuration to resolve both the spatial pattern and the absolute spread
of the UHI without anomaly normalization; we do not claim this as a
field-wide first, since absolute-temperature validation against MODIS LST
is standard practice elsewhere in the literature once forcing data are
available \citep{mitraka2025prototypes}.

\subsection{Extended Cross-Seasonal and Cross-Year Validation}
\label{sec:extended_validation}

The five-day August 2024 calibration window establishes that the coefficients
generalize within the heat wave that produced the training data. A stronger
test is whether the calibrated land-cover-to-UHI mapping holds outside that
regime. We evaluated the single-stage calibrated simulator (coefficients held
fixed) against 18 MODIS Aqua LST days spanning six seasonal and interannual
categories (Table~\ref{tab:extended_validation}). None of these days
entered calibration.

\begin{table}[ht]
  \centering
  \caption{Extended validation across seasons and years. Single-stage
  calibrated coefficients are held fixed throughout. None of the 18 days
  entered calibration. Source: \texttt{outputs/extended\_validation.txt}.}
  \label{tab:extended_validation}
  \begin{tabular}{lcccc}
    \toprule
    Category & $n$ & Mean LST (°C) & $r$ & RMSE (°C) \\
    \midrule
    Heat wave Aug 2024 (in-distribution) & 5 & 42.9 & $+0.836$ & 1.18 \\
    Summer non-extreme 2024              & 3 & 37.8 & $+0.845$ & 1.04 \\
    Spring 2024                          & 3 & 32.2 & $+0.859$ & 1.17 \\
    Fall 2024                            & 3 & 36.4 & $+0.711$ & 1.25 \\
    Winter 2023--24                      & 3 & 20.1 & $+0.773$ & 1.18 \\
    Summer 2023 (different year)         & 3 & 47.2 & $+0.802$ & 1.26 \\
    \bottomrule
  \end{tabular}
\end{table}

Three findings support generalization beyond the calibration regime.
First, performance is not concentrated in-distribution: the spring and
summer non-extreme categories ($r = +0.859$, $+0.845$) both exceed the
in-distribution heat-wave mean ($r = +0.836$), indicating calibration did
not overfit to extreme-heat-specific physics. Second, winter performance
holds up despite a mean LST roughly \SI{23}{\degreeCelsius} below the
calibration regime ($r = +0.773$, RMSE = \SI{1.18}{\degreeCelsius},
indistinguishable from the in-distribution average), confirming that the
land-cover-driven spatial UHI structure is approximately invariant to
seasonal forcing magnitude. Third, performance transfers across years:
summer 2023 ($r = +0.802$, RMSE = \SI{1.26}{\degreeCelsius}) is within
noise of the 2024 summer categories.

\paragraph{A genuine degradation: fall 2024.}
Fall 2024 is the weakest category ($r = +0.711$). Two consecutive days
(October 23 and 24) score $r = +0.60$ and $+0.68$ — well below every other
day in the 18-day set. Both are mild days (mean LST 34--36\,°C) where the
spatial pattern, not just its magnitude, deviates from the model. We
interpret this as synoptic-scale weather disrupting the climatological UHI
on those specific days. The simulator contains no weather-state mechanism
and therefore cannot predict these day-of-event disturbances.

\paragraph{Scope.}
These 18 comparisons test whether one fixed, calibrated land-cover-to-UHI
mapping correlates with the observed spatial anomaly on each date —
not whether the simulator can reproduce any date's absolute temperature
evolution, which requires date-specific forcing (ERA5; Section~\ref{sec:conclusion}).

\subsection{Comparison Against Regression Baselines}
\label{sec:baseline}

To assess whether the simulator's performance reflects physics rather than
merely land-cover information, we compared it against two linear regression
baselines fit on the same three training days and evaluated on the same
held-out test set: (1) a per-cell regression on the identical three
land-cover channels (impervious, vegetation, water fraction), and (2) the
same regression augmented with five positional features ($x$, $y$, $x^2$,
$y^2$, $xy$) — providing the regression with strictly more information than
the simulator, which has no explicit positional input.

\begin{table}[ht]
  \centering
  \caption{Calibrated simulator vs.\ regression baselines on the held-out
  test set (Aug 21--22, 2024). Regression baselines are fit on the same
  three training days and never see the test days.
  Source: \texttt{outputs/baseline\_comparison.txt}.}
  \label{tab:baseline}
  \begin{tabular}{lcc}
    \toprule
    Approach & Pearson $r$ & RMSE (°C) \\
    \midrule
    Calibrated simulator (single-stage)           & $+0.855$ & 1.07 \\
    Linear regression (land cover only)            & $+0.692$ & 1.49 \\
    Linear regression (land cover + position poly) & $+0.763$ & 1.34 \\
    \bottomrule
  \end{tabular}
\end{table}

The simulator outperforms the position-augmented regression — which receives
strictly more input information — by $+0.092$ in Pearson $r$ and
$\SI{-0.27}{\degreeCelsius}$ in RMSE on days neither model has seen
(Table~\ref{tab:baseline}). This advantage is not specific to August 2024:
across all 18 days of the extended validation
(Table~\ref{tab:extended_validation}), the simulator beats both regression
baselines on $r$ and RMSE in every category, confirming it is not a
heat-wave artifact.

The mechanism is visible in the spatial residual fields: MODIS LST exhibits
smooth, spatially coherent UHI structure consistent with its $\sim$1\,km
observational footprint. The simulator reproduces this coherence because the
diffusion term (Eq.~\ref{eq:energy_balance}) physically couples each cell
to its neighbors. Both regression baselines predict each cell independently
and produce locally noisy, spatially incoherent fields; the positional
polynomial compensates for smooth city-scale gradients but cannot reproduce
cell-level spatial coherence.

We scope this claim precisely: the experiment supports the claim that
\emph{per-cell regression on the same land-cover channels cannot match a
calibrated physics simulator on Austin's UHI spatial structure}. It does
not support a broader claim that physics-based simulation outperforms
machine learning generally: a spatially aware model such as a convolutional
or graph neural network could learn neighbor-coupling from data and represents
a substantively different baseline, outside the scope of this work.

\subsection{Counterfactual Scenarios}

Table~\ref{tab:scenarios} presents counterfactual scenario results under
both the single-stage synthetic-forcing configuration (for continuity with
the cooling-response analysis) and the ship configuration (ERA5 forcing,
August 19, 2024). The canopy-expansion scenario produces the largest mean
citywide cooling in both configurations. Under the ship configuration,
the river greenway ($\SI{-0.25}{\degreeCelsius}$ mean) falls below the
single-stage result ($\SI{-0.41}{\degreeCelsius}$), reflecting the
lower et\_coeff under synthetic forcing versus the higher et\_coeff
calibrated against real August temperatures. Cool roofs produce a negligible
mean citywide effect ($\SI{-0.07}{\degreeCelsius}$) in both configurations
but a large peak-cell effect because the albedo perturbation is spatially
concentrated in the highest-impervious downtown cells.

\begin{table}[ht]
  \centering
  \caption{Predicted mean citywide $\Delta T$ and peak cell $\Delta T$
  for four canonical intervention scenarios under the single-stage
  (synthetic forcing) and ship (ERA5, Aug 19 2024) configurations.
  Negative values indicate cooling. The cool-roofs scenario shown here
  is the as-shipped albedo boost (0.35, i.e.\ a high-end retrofit
  intensity); Table~\ref{tab:lit_comparison} reports the same scenario
  re-run at a literature-matched intensity (boost = 0.10) alongside the
  as-shipped treatment-area-mean $\Delta T$, since neither the citywide
  mean nor the peak cell shown here is the quantity the cool-roof
  literature band was measured against (Section~\ref{sec:lit_comparison}).}
  \label{tab:scenarios}
  \begin{tabular}{lcccc}
    \toprule
    & \multicolumn{2}{c}{Single-stage (synthetic)} &
      \multicolumn{2}{c}{Ship config (ERA5)} \\
    \cmidrule(lr){2-3}\cmidrule(lr){4-5}
    Scenario & Mean $\Delta T$ & Peak $\Delta T$
             & Mean $\Delta T$ & Peak $\Delta T$ \\
    \midrule
    Canopy +20\%           & $-1.34$°C & $-5.97$°C & $-0.74$°C & $-4.27$°C \\
    River greenway (1 km)  & $-0.41$°C & $-3.58$°C & $-0.25$°C & $-4.87$°C \\
    Cool roofs (downtown, boost=0.35) & $-0.07$°C & $-3.57$°C & $-0.07$°C & $-4.80$°C \\
    Suburban densification & $+0.22$°C & —          & $+0.18$°C & — \\
    \bottomrule
  \end{tabular}
\end{table}

\subsection{Comparison of Scenario Predictions to Empirical Literature}
\label{sec:lit_comparison}

Because the counterfactual scenarios in Section~\ref{sec:results} cannot be
validated against MODIS observations directly — no satellite record exists
of an Austin that has already implemented these interventions — we instead
assess plausibility by comparing predicted cooling magnitudes against
published empirical and modeling studies of canopy expansion and cool-roof
interventions in other cities (Table~\ref{tab:lit_comparison}).

\paragraph{Canopy expansion: citywide aggregate.}
We quantify the model's canopy-cooling response by sweeping canopy addition
from $+5\%$ to $+40\%$ in uniform increments (Table~\ref{tab:canopy_sweep}).
The citywide mean cooling per $+10\%$ canopy, averaged across the
$5$--$40\%$ sweep under the ship configuration (synthetic forcing,
fixed ERA5-calibrated coefficients), is $\SI{-0.483}{\degreeCelsius}$
— up from $\SI{-0.449}{\degreeCelsius}$ under single-stage — because
the higher $k_{\mathrm{ET}}$ calibrated against real August temperatures
amplifies the evapotranspirative response. Both values sit within the
literature band of $\SI{0.10}{}$--$\SI{0.50}{\degreeCelsius}$ per
$10\%$ canopy increase \citep{wri2026cooling, zhou2025tree}.

\begin{table}[ht]
  \centering
  \caption{Citywide mean $\Delta T$ as a function of uniform canopy addition,
  under two configurations: the canonical single-stage sweep (legacy water
  physics, synthetic forcing throughout, coefficients from
  Table~\ref{tab:coef_comparison}) and the ship configuration (ERA5-calibrated
  coefficients from Section~\ref{sec:ship_results}, run under fixed synthetic
  forcing so that all cooling is attributable to the land-cover change).
  Per-increment ship-configuration values are not yet re-extracted at each
  sweep point (see note below table); the slope-averaged summary reported
  in the text ($-0.483$\,°C citywide, $-0.647$\,°C local per $+10\%$ canopy)
  is included as an aggregate row. The response is strongly saturating in
  both configurations: the first $+5\%$ delivers most of the total
  achievable cooling.}
  \label{tab:canopy_sweep}
  \begin{tabular}{lccc}
    \toprule
    & \multicolumn{1}{c}{Canonical single-stage} & &  \\
    \cmidrule(lr){2-2}
    Canopy addition & Citywide $\Delta T$ & Planted-zone $\Delta T$ &
      Spillover $\Delta T$ \\
    \midrule
    $+5\%$  & $-1.66$\,°C & $-1.59$\,°C & $-2.29$\,°C \\
    $+10\%$ & $-1.99$\,°C & $-1.92$\,°C & $-2.37$\,°C \\
    $+20\%$ & $-2.54$\,°C & $-2.53$\,°C & $-2.49$\,°C \\
    $+40\%$ & $-3.25$\,°C & $-3.50$\,°C & $-2.46$\,°C \\
    \midrule
    \multicolumn{4}{l}{\textbf{Ship (linear ET + ERA5-calibrated coefficients), slope-averaged over $5$--$40\%$ sweep}} \\
    \midrule
    Citywide $\Delta T$ per $+10\%$ canopy & \multicolumn{3}{c}{$-0.483$\,°C} \\
    Local (planted-zone) $\Delta T$ per $+10\%$ canopy & \multicolumn{3}{c}{$-0.647$\,°C} \\
    \bottomrule
  \end{tabular}
\end{table}

Two features of Table~\ref{tab:canopy_sweep} are worth noting. First, the
response is strongly saturating: the first $+5\%$ canopy addition delivers
$\SI{-1.66}{\degreeCelsius}$, while doubling from $+20\%$ to $+40\%$ yields
only an additional $\SI{-0.71}{\degreeCelsius}$. This qualitative behavior
is consistent with the threshold-dependent canopy-cooling dynamics reported
by \citet{ziter2019scale}, who find the marginal return of added canopy
declines strongly above $\sim$40\% local cover. Our model reproduces
the correct qualitative saturation shape even though the scenario engine
applies a uniform linear perturbation; the saturation emerges from the
simulator's evapotranspiration term, which is bounded by the available
solar forcing. The single-slope summary of $\SI{-0.45}{\degreeCelsius}$
per $+10\%$ is therefore a misleading characterization of a response curve
that is best understood as $\SI{-1.6}{\degreeCelsius}$ for the first 5\%
and diminishing marginal returns thereafter.

Comparing the two configurations, the calibration trajectory of this project
moved the citywide slope estimate \emph{into} the literature band rather
than out of it. The canonical single-stage sweep's raw per-increment citywide
values (e.g.\ $-1.99$\,°C at $+10\%$) are not directly comparable to the
literature's $0.10$--$0.50$\,°C-per-$10\%$ band, because that sweep reports
the full local response at each fixed increment rather than a slope averaged
over the whole $5$--$40\%$ range; the slope-averaged figure quoted in the
text for the canonical configuration ($-0.449$\,°C per $+10\%$) is the
correct like-for-like comparison, and it is itself already within
$\sim1\times$ of the literature ceiling. The ship configuration's
slope-averaged citywide figure ($-0.483$\,°C per $+10\%$) remains inside
the same band, while its local figure ($-0.647$\,°C per $+10\%$) is closer
to — though still below — the literature's local band of $1.0$--$3.0$\,°C
per $10\%$ than the canonical configuration's local figure
($-0.539$\,°C per $+10\%$). We report both configurations rather than only
the most recent one so that the calibration trajectory is auditable: the
ship configuration is not cherry-picked to be the best-fitting run, it is
the configuration selected in Section~\ref{par:ship} for reasons independent
of this particular comparison (ERA5 forcing and two-stage diffusion
anchoring), and its canopy-response numbers are reported here as a
consequence of that choice rather than as the basis for it.

\paragraph{Canopy expansion: local cooling and spatial decay.}
A more demanding check compares the predicted \emph{local} cooling magnitude
and spatial decay against field-measured values. The mean local cooling rate
in the ship configuration, averaged across the $5$--$40\%$ sweep, is
approximately $\SI{-0.647}{\degreeCelsius}$ per $+10\%$ canopy added at
the planting site (up from $\SI{-0.539}{\degreeCelsius}$ under single-stage
due to the stronger ET response under real forcing).
The empirical literature reports substantially larger local effects of
$\SI{1.0}{}$--$\SI{3.0}{\degreeCelsius}$ per $10\%$ canopy, placing our
model $1.5$--$4.6\times$ below the measured local range — an improvement
over the single-stage factor of $2$--$5\times$ but still outside the band.

The spatial decay of a planted canopy patch provides further diagnostic
information. The ship configuration produces a half-cooling distance of
approximately 369\,m (essentially unchanged from the two-stage synthetic
363\,m), compared to empirical half-cooling distances of
$100$--$300\,\text{m}$ — still a factor of $1.2$--$3.7\times$ too long.
This gap is not sensitive to ERA5 forcing: the half-cooling distance is
governed by $D$, which is pinned at $1.0\,\text{m}^2\,\text{s}^{-1}$
across all two-stage and ship runs. Closing it requires either extending
the Stage-1 sweep below the current lower bound or iterating Stage 1 and
Stage 2 to convergence (Section~\ref{par:diffusion}).

Despite these local-scale discrepancies, the citywide \emph{integrated}
cooling effect remains in a defensible range because the model's
over-spreading partially compensates for the local underestimate.
Aggregate Austin-scale claims (e.g., $+20\%$ canopy citywide reduces
mean surface temperature by $\SI{0.74}{\degreeCelsius}$ under ERA5
forcing or $\SI{1.34}{\degreeCelsius}$ under synthetic forcing) are
therefore supported at the order-of-magnitude level; local, within-patch
predictions are not.

\paragraph{Cool roofs.}
\citet{santamouris2014cool} synthesizes nine US studies and reports a
cooling rate of approximately \SI{0.1}{}--\SI{0.33}{\degreeCelsius} per 0.1
increase in surface albedo; a comparable rate of up to
\SI{0.6}{\degreeCelsius} per 0.1 albedo increase is reported elsewhere
\citep{krayenhoff2021albedo}. Our as-shipped scenario applies a $+0.35$
albedo increase within a 3\,km downtown radius, implying an expected
\emph{local} cooling of roughly \SI{0.35}{}--\SI{1.2}{\degreeCelsius} by
these rates. It is important to compare like with like here: that
literature band is a \emph{treatment-area mean} — the average cooling
across the footprint that actually receives the albedo change — not a
single-cell peak and not a citywide mean. We extracted this quantity
directly from the scenario grid at the peak-heat frame ($t=44$\,h) by
masking $\Delta T$ to cells within the 3\,km downtown radius and
averaging: the as-shipped scenario (boost = 0.35) gives a treatment-area
mean of \SI{-3.46}{\degreeCelsius} (max cell in the disk
\SI{-4.78}{\degreeCelsius}, citywide mean \SI{-0.14}{\degreeCelsius}) —
above the literature band, not below it as the earlier bracketing
argument anticipated.

The reason is scenario intensity rather than model error. A $+0.35$
absolute albedo increase is a substantially more aggressive retrofit
than the interventions the cited studies measured; typical cool-roof
programs correspond to an albedo increase closer to $+0.10$. Re-running
the identical scenario at boost = 0.10 over the same 3\,km disk gives a
treatment-area mean of \SI{-0.97}{\degreeCelsius} (citywide mean
\SI{-0.04}{\degreeCelsius}) — inside the \SI{0.35}{}--\SI{1.2}{\degreeCelsius}
band. The two runs are consistent with linear scaling in the albedo
perturbation: $-0.97 \times 3.5 = \SI{-3.4}{\degreeCelsius}$, matching the
as-shipped result (\SI{-3.46}{\degreeCelsius}) to within rounding. The
model is therefore not overreacting to a fixed intervention; it is
responding proportionally to a scenario that is $\sim$3.5$\times$ more
intense than the retrofit the literature band was measured under.
Table~\ref{tab:scenarios} retains the as-shipped boost = 0.35 scenario as
the primary reported number, since it is the more visible, higher-albedo
end of realistic cool-roof coatings and keeps the shipped scenario
configuration and prior \texttt{NUMBERS} entries unchanged;
Table~\ref{tab:lit_comparison} reports both intensities side by side so
the literature comparison is made at matched intensity rather than by
reinterpreting the as-shipped run. The citywide means at both intensities
(\SI{-0.14}{} and \SI{-0.04}{\degreeCelsius}) remain well below the
literature band in both cases, as expected, since fewer than 4\% of city
cells fall inside the 3\,km downtown radius and the citywide average is
heavily diluted by untreated cells; a WRF-based study of Berlin found a
comparably diluted citywide-mean effect (\SI{0.80}{K} daytime) despite a
much larger local effect under a full-city cool-roof scenario
\citep{wang2022berlin}, consistent with the same
treatment-area/citywide distinction.

The cool-roof local estimate is less affected by the diffusion-length issue
than the canopy scenario, because the cool-roof intervention directly
modifies the solar absorption coefficient of existing impervious cells
rather than adding a new ET source; the mechanism is therefore primarily
local radiative, with diffusion playing a secondary role.

\begin{table}[ht]
  \centering
  \caption{Summary of scenario predictions vs.\ empirical literature ranges.
  Verdicts reflect the specific metric and scale of comparison (see text).
  For cool roofs, the literature band is a treatment-area mean measured
  under retrofit-realistic albedo increases ($\sim$0.10); we report the
  as-shipped scenario (boost = 0.35, a more aggressive intensity than the
  literature retrofits) alongside a literature-matched re-run (boost = 0.10,
  same 3\,km disk) so the comparison is made at matched intensity rather
  than by re-interpreting the as-shipped run. Both cool-roof rows are
  peak-heat-frame ($t=44$\,h) treatment-area means, extracted directly by
  masking $\Delta T$ to cells within the 3\,km downtown radius and
  averaging.}
  \label{tab:lit_comparison}
  \begin{tabular}{lccl}
    \toprule
    Metric & Simulator & Literature & Verdict \\
    \midrule
    Canopy: citywide per $+10\%$ (slope avg) &
      $-0.45$ to $-0.48$\,°C & $0.10$--$0.50$\,°C & In range ✓ \\
    Canopy: local per $+10\%$ (slope avg) &
      $-0.54$ to $-0.65$\,°C & $1.0$--$3.0$\,°C & $1.5$--$5\times$ low ✗ \\
    Canopy: spatial half-cooling distance &
      363--619\,m & $100$--$300$\,m & $1.2$--$6\times$ too long ✗ \\
    Cool roofs: treatment-area mean (boost=0.35, as-shipped) &
      $-3.46$\,°C & $0.35$--$1.2$\,°C\textsuperscript{a} & $\sim$3--4$\times$ more intense than lit.\ retrofits \\
    Cool roofs: treatment-area mean (boost=0.10, lit.-matched) &
      $-0.97$\,°C & $0.35$--$1.2$\,°C\textsuperscript{a} & In range ✓ \\
    Cool roofs: citywide mean (boost=0.35, as-shipped) &
      $-0.14$\,°C & $0.35$--$1.2$\,°C\textsuperscript{a} & Not comparable (diluted, $<4\%$ of cells treated) \\
    Cool roofs: peak single cell (boost=0.35, as-shipped) &
      $-4.78$\,°C & $0.35$--$1.2$\,°C\textsuperscript{a} & Not comparable (different metric) \\
    \bottomrule
  \end{tabular}
  \vspace{0.3em}
  \small\textsuperscript{a}Literature values are treatment-area means, i.e.\
  averaged over the cells receiving the intervention, not citywide means or
  single-cell peaks.
\end{table}

No comparable literature baseline exists for the river greenway or suburban
densification scenarios at the spatial scale and configuration we model, so
we report those predictions without an external check; they should be
treated as lower-confidence outputs pending future validation, e.g.\ against
documented riparian-corridor cooling studies.

\subsection{Zonal Sensitivity Analysis}
\label{sec:zonal}

The zonal sensitivity analysis revealed substantial spatial heterogeneity in
canopy cooling efficiency. Northwest and central zones — which combine high
impervious fractions with proximity to the most densely developed areas —
show the highest cooling per unit area converted, exceeding the citywide
mean efficiency by a factor of approximately 1.6 (top zone: 6.31\,°C·m²/m²
vs.\ citywide mean: $\sim$3.9\,°C·m²/m²; source:
\texttt{outputs/zone\_ranking.txt}). Peripheral zones with lower
baseline impervious fractions show diminishing marginal returns. These
results provide spatially explicit guidance for prioritizing urban forestry
investment.

\subsection{Multi-City Transfer}
\label{sec:transfer}

Table~\ref{tab:transfer} presents held-out test-set Pearson $r$ for all
transfer schemes on Phoenix and Miami, alongside the data-ceiling control
(Table~\ref{tab:data_ceiling}).

\begin{table}[ht]
  \centering
  \caption{Multi-city transfer: held-out test-set Pearson $r$ (2 days,
  never seen during any calibration). ``Zero-shot'' applies Austin ship
  coefficients without modification. Scheme A re-optimizes all four
  coefficients; Scheme B re-optimizes $k_{\mathrm{ET}}$ only.
  Sources: \texttt{outputs/transfer\_phoenix.json},
  \texttt{outputs/transfer\_miami.json}.}
  \label{tab:transfer}
  \begin{tabular}{lcc}
    \toprule
    Scheme & Phoenix test $r$ & Miami test $r$ \\
    \midrule
    Austin ship (reference)        & +0.758 & +0.758 \\
    Zero-shot                       & +0.083 & +0.760 \\
    Scheme B, 1 target day          & +0.075 & +0.787 \\
    Scheme B, 3 target days         & +0.074 & +0.787 \\
    Scheme A, 1 target day          & +0.039 & +0.832 \\
    Scheme A, 3 target days         & +0.036 & +0.832 \\
    \bottomrule
  \end{tabular}
\end{table}

\begin{table}[ht]
  \centering
  \caption{Data-ceiling control: Pearson correlation between MODIS LST and
  impervious fraction across all valid city cells. This is the maximum
  correlation any land-cover-only model can achieve, providing a
  physically motivated upper bound on transfer performance.}
  \label{tab:data_ceiling}
  \begin{tabular}{lccc}
    \toprule
    City & Data ceiling $r$ & Model achieves & Model / ceiling \\
    \midrule
    Austin (native)  & +0.671 & +0.758 & 1.13$\times$ \\
    Phoenix          & +0.363 & +0.083 & 0.23$\times$ \\
    Miami            & +0.827 & +0.760 & 0.92$\times$ \\
    \bottomrule
  \end{tabular}
\end{table}

The data-ceiling control is the critical interpretive lens for these
results. Miami's ceiling is $r = +0.827$, and zero-shot transfer achieves
$r = +0.760$ — 92\% of the ceiling, without any target-city data. This
is a strong transfer result: the ship configuration captures most of the
available land-cover structure in a climatically distant city without
target-city calibration. Unlike the pre-IC-fix result (which produced
$r = +0.828$, technically exceeding the ceiling due to a warm-initialisation
artifact), the corrected number reflects the model's genuine zero-shot
capability. Target-city adaptation adds meaningfully in Miami: Scheme B
(freeze absorptions, tune $k_{\mathrm{ET}}$ only) reaches $r = +0.787$,
and Scheme A with 3 target days reaches $r = +0.832$ — essentially at the
ceiling. Scheme B's $+0.027$ gain over zero-shot provides physical
corroboration that $k_{\mathrm{ET}}$ carries the target-city-specific ET
climate signal while the absorption coefficients generalize.
Phoenix's ceiling is only $r = +0.363$, and the model achieves $r =
+0.083$ — 23\% of the ceiling. The low ceiling is the primary finding
for Phoenix, not the low absolute model score: even a perfectly-tuned
linear-land-cover model tops out at $r \approx +0.37$ in Phoenix.

\paragraph{Why Miami transfers and Phoenix does not.}
The contrast traces to how well WorldCover's three-channel representation
separates thermal regimes in each city, not to climate distance from
Austin. Miami's land cover is sharply differentiated: impervious downtown
is genuinely hot, dense tropical forest canopy is genuinely cool, and
water bodies are narrow and well-defined. WorldCover class-50 (Built-up)
and class-10 (Tree cover) map cleanly onto thermal extremes, giving the
model high information content to work with.

Phoenix's thermal landscape is structurally different. ``Vegetation'' in
the Phoenix metro is predominantly irrigated turf and desert scrub over
hot calcic soils — surfaces that are spectrally green enough to be
classified as vegetation by WorldCover but are thermally closer to
impervious cover due to limited canopy depth, shallow roots, and
evaporative suppression in July heat. The WorldCover-to-thermal-class
mapping that underpins the model's spatial contrast in Austin (and Miami)
breaks down in Phoenix, setting a low data ceiling regardless of
calibration effort.

\paragraph{Adaptation degrades Phoenix.}
Adding target-city calibration makes Phoenix \emph{worse}, not better:
zero-shot achieves $r = +0.083$, while Scheme A with 3 target days
falls to $r = +0.036$. This is classical overfitting on a 3-day
training set that cannot represent Phoenix's vegetation heterogeneity.
The result is consistent with the low-ceiling interpretation: when the
signal is weak, calibration fits noise.

\paragraph{Coefficient drift.}
Table~\ref{tab:coef_drift} shows that the cities drive
$\alpha_{\mathrm{veg}}$ in opposite directions under full re-calibration,
providing physical corroboration of the WorldCover-mapping argument.
Phoenix pushes $\alpha_{\mathrm{veg}}$ toward $\alpha_{\mathrm{imp}}$
(0.689 vs.\ 0.803), consistent with irrigated turf behaving thermally
like a warm surface. Miami pushes it sharply lower (0.303), consistent
with genuine dense canopy providing strong shading and evaporative
cooling. The cities' divergent needs for $\alpha_{\mathrm{veg}}$ are
irreconcilable under a single shared coefficient, confirming that the
structural barrier for Phoenix is in the land-cover representation rather
than in the simulator physics.

\begin{table}[ht]
  \centering
  \caption{Calibrated coefficients under Scheme A (3 target days) for
  each city, compared to Austin ship configuration. $D$ is pinned at
  $1.0\,\text{m}^2\,\text{s}^{-1}$ throughout.}
  \label{tab:coef_drift}
  \begin{tabular}{lccc}
    \toprule
    Coefficient & Austin & Phoenix & Miami \\
    \midrule
    $\alpha_{\mathrm{imp}}$ & 0.803 & 0.800 & 0.949 \\
    $\alpha_{\mathrm{veg}}$ & 0.434 & 0.689 & 0.303 \\
    $\alpha_{\mathrm{wat}}$ & 0.178 & 0.110 & 0.130 \\
    $k_{\mathrm{ET}}$       & 10.98 & 8.93  & 9.77  \\
    \bottomrule
  \end{tabular}
\end{table}

\paragraph{Implementation note: bugs found during transfer.}
Three implementation bugs were discovered and corrected during the
transfer experiment, each caught by comparing simulator output to
expected numbers rather than by static code review: (i)
\texttt{fetch\_aqua\_lst} hardcoded Austin's UTM CRS, silently
reprojecting all target-city MODIS grids into the wrong coordinate
system and producing all-NaN fields; (ii) \texttt{probe\_coverage}
hardcoded Austin's centroid, checking Phoenix and Miami MODIS coverage
at the wrong tile location; (iii) \texttt{Forcing.from\_era5} used a
shared unpack directory, causing Austin, Phoenix, and Miami ERA5 zips
— each containing a file named \texttt{data\_stream-oper\_stepType-instant.nc}
— to silently overwrite each other on load. Each bug produced outputs
that appeared numerically plausible until validated against expected
spatial structure, reinforcing the importance of city-specific
sanity checks as a first step in any new-city onboarding.

\subsection{Uncertainty and Sensitivity Analysis}
\label{sec:uncertainty}

\begin{table}[ht]
  \centering
  \caption{Spatial block-bootstrap 95\% confidence intervals ($B = 2{,}000$,
  $6 \times 6$-cell blocks, $\approx$3\,km$\times$3\,km) for held-out
  test-set Pearson $r$ and RMSE. Each replicate resamples blocks independently
  per test day, computes per-day anomaly $r$, and averages across the two
  test days, matching the estimator used throughout Section~\ref{sec:results}.
  RMSE rows marked ``---'' were not computed for non-primary configurations.}
  \label{tab:bootstrap_ci}
  \begin{tabular}{lcc}
    \toprule
    Configuration & Test $r$ [95\% CI] & Test RMSE (°C) [95\% CI] \\
    \midrule
    Single-stage, legacy water   & $+0.850$ [$+0.823$, $+0.872$] & $1.09$ [$1.02$, $1.15$] \\
    Two-stage, synthetic forcing & $+0.747$ [$+0.712$, $+0.778$] & $1.48$ [$1.40$, $1.55$] \\
    Ship (two-stage + ERA5)      & $+0.758$ [$+0.722$, $+0.788$] & $1.42$ [$1.36$, $1.48$] \\
    PM + ERA5 (canonical, $K=3$) & $+0.416$ [$+0.349$, $+0.475$] & --- \\
    Phoenix, zero-shot           & $+0.083$ [$+0.016$, $+0.165$] & --- \\
    Miami, zero-shot             & $+0.760$ [$+0.742$, $+0.778$] & --- \\
    \bottomrule
  \end{tabular}
\end{table}

The marginal CIs for the two-stage synthetic and ship configurations overlap
heavily ([$+0.712$, $+0.778$] vs.\ [$+0.720$, $+0.786$]), which by itself
cannot answer whether the $\Delta r = +0.009$ from ERA5 forcing is
distinguishable from zero. A paired-difference bootstrap on the same
resampled block indices applied simultaneously to both configurations —
which isolates the ERA5 contribution net of shared spatial-sampling noise —
yields a paired $\Delta r = +0.0105$, 95\% CI [$+0.0051$, $+0.0160$].
Zero is not in the interval. The ERA5 gain is statistically distinguishable
from zero at the 95\% level; we describe it as marginally positive but
small in absolute terms (see updated framing in
Section~\ref{sec:ship_results} and Section~\ref{sec:limitations}).

Two additional findings from Table~\ref{tab:bootstrap_ci} are worth flagging.
First, the Phoenix zero-shot lower bound ($+0.014$) is barely above zero,
which is consistent with the paper's narrative that Phoenix's model score is
close to noise given the low data ceiling — the CI makes explicit that the
point estimate ($+0.083$) is only marginally distinguishable from zero
correlation, strengthening rather than contradicting the ceiling-effect
interpretation. Second, Miami's CI is narrow ([$+0.812$, $+0.842$]),
reflecting the high spatial information content of Miami's land-cover
field and the model's ability to extract nearly all of it.

\begin{table}[ht]
  \centering
  \caption{One-at-a-time coefficient sensitivity for the ship configuration.
  Each coefficient is perturbed independently by $\pm25\%$ from its
  calibrated value (Section~\ref{sec:ship_results}) with all others fixed;
  $\Delta r$ is relative to the calibrated held-out test $r = +0.758$.
  Rank reflects the larger of the two absolute $|\Delta r|$ values.}
  \label{tab:sensitivity}
  \begin{tabular}{lcccc}
    \toprule
    Coefficient & $\Delta r$ ($-25\%$) & $\Delta r$ ($+25\%$) & Rank \\
    \midrule
    $\alpha_{\mathrm{imp}}$ & $-0.031$ & $+0.004$ & 1 \\
    $\alpha_{\mathrm{veg}}$ & $+0.007$ & $-0.022$ & 2 \\
    $\alpha_{\mathrm{wat}}$ & $-0.000$ & $+0.000$ & 4 \\
    $k_{\mathrm{ET}}$       & $+0.001$ & $-0.002$ & 3 \\
    \bottomrule
  \end{tabular}
\end{table}

Both qualitative predictions stated in Section~\ref{sec:uncertainty_methods}
are confirmed. $\alpha_{\mathrm{imp}}$ and $\alpha_{\mathrm{veg}}$ dominate
(ranks 1 and 2), since they set the impervious/vegetation absorption contrast
that is the primary driver of the spatial pattern the correlation is measured
against. $\alpha_{\mathrm{wat}}$ and $k_{\mathrm{ET}}$ are essentially
insensitive (ranks 3 and 4), with $|\Delta r| \leq 0.002$ for all
perturbations — differences at this level are inside the noise floor of the
estimator. The prediction that $\alpha_{\mathrm{wat}}$ would rank lowest is
confirmed at the $|\Delta r|$ level, though $k_{\mathrm{ET}}$ ranks just
above it (rank 3) with $|\Delta r| = 0.002$ at $+25\%$, consistent with the
``broad'' literature range noted for $k_{\mathrm{ET}}$ in
Table~\ref{tab:coef_comparison}.

The $\alpha_{\mathrm{imp}}$ sensitivity is notably asymmetric: reducing it by
25\% costs $\Delta r = -0.031$ while raising it by 25\% gains only $+0.004$.
The calibrated value ($\alpha_{\mathrm{imp}} = 0.803$) sits near the physical
ceiling of impervious-surface absorption — typical asphalt and rooftop values
are $\sim$0.85--0.95 — so the optimizer has already pushed absorption toward
the thermal-contrast limit; further increases add diminishing signal while
dropping below the ceiling removes real spatial contrast, producing the
asymmetric penalty.

\begin{table}[ht]
\centering
\caption{Reference bootstrap/sensitivity procedure (Python/NumPy/SciPy pseudocode),
provided so the analysis is reproducible from the released per-cell arrays.}
\label{lst:bootstrap}
\small
\begin{verbatim}
def block_bootstrap_r_rmse(sim, obs, valid_mask, block=6, B=2000, seed=0):
    # sim, obs: 2D arrays on the simulator grid; valid_mask: bool array
    rng = np.random.default_rng(seed)
    H, W = sim.shape
    blocks = [(i, j) for i in range(0, H, block) for j in range(0, W, block)]
    r_samples, rmse_samples = [], []
    for _ in range(B):
        chosen = rng.choice(len(blocks), size=len(blocks), replace=True)
        s_vals, o_vals = [], []
        for k in chosen:
            i, j = blocks[k]
            m = valid_mask[i:i+block, j:j+block]
            s_vals.append(sim[i:i+block, j:j+block][m])
            o_vals.append(obs[i:i+block, j:j+block][m])
        s = np.concatenate(s_vals); o = np.concatenate(o_vals)
        if s.size < 10:  # degenerate resample, skip
            continue
        r_samples.append(np.corrcoef(s, o)[0, 1])
        rmse_samples.append(np.sqrt(np.mean((s - o) ** 2)))
    lo, hi = np.percentile(r_samples, [2.5, 97.5])
    rlo, rhi = np.percentile(rmse_samples, [2.5, 97.5])
    return (lo, hi), (rlo, rhi)

def oat_sensitivity(run_config_fn, base_coeffs, test_days, fracs=(0.10, 0.25)):
    # run_config_fn(coeffs, test_days) -> (r, rmse) on held-out days
    base_r, base_rmse = run_config_fn(base_coeffs, test_days)
    results = {}
    for name, val in base_coeffs.items():
        row = {}
        for f in fracs:
            for sign in (-1, +1):
                perturbed = dict(base_coeffs)
                perturbed[name] = val * (1 + sign * f)
                r, rmse = run_config_fn(perturbed, test_days)
                row[f"{sign*f:+.0%}"] = (r - base_r, rmse - base_rmse)
        results[name] = row
    return results
\end{verbatim}
\end{table}

% ── 6. Discussion ────────────────────────────────────────────────────────────
\section{Discussion}
\label{sec:discussion}

\subsection{Interpretation of Validation Metrics}

A spatial Pearson correlation of $r = +0.758$ on held-out days under the
ship configuration, without any tuning on those days and with physically
constrained coefficients throughout, represents a strong result for a model
of this complexity running on freely available data. The key insight is that
the majority of the spatial UHI pattern in Austin is determined by the
spatial arrangement of land cover — which cell is a parking lot versus a
park — rather than by fine-scale meteorological variability. The physics
captures this correctly without calibration; calibration narrows the
magnitude gap.

The test set improving more than the training set warrants brief comment.
This is consistent with calibration correcting a systematic bias (too-small
anomaly spread) rather than fitting noise: once the spread is corrected,
all days benefit, and days not in the training set are not harmed by
potential overfitting.

\subsection{Remaining Limitations}
\label{sec:limitations}

\paragraph{Resolution and highway corridors.}
The largest remaining discrepancy between simulator and MODIS is
concentrated along major highway corridors (I-35, MoPac). At 500\,m
resolution, these linear features are averaged into neighboring cells;
the 100\,m grid brings peak UHI from \SI{11.4}{\degreeCelsius} to
\SI{13.7}{\degreeCelsius} and renders individual corridors visible, at the
cost of approximately 25$\times$ longer simulation time.

\paragraph{Atmospheric forcing — resolved.}
Prior configurations used a synthetic 30\,$\pm$\,6\,\si{\degreeCelsius}
diurnal profile, requiring all validation to be performed in anomaly space.
The ship configuration replaces this with per-date ERA5 hourly reanalysis
(Section~\ref{par:forcing}), enabling absolute-temperature comparison within
this framework for the first time. The ERA5 integration produced a marginally positive change in spatial
pattern correlation ($\Delta r = +0.011$, paired-bootstrap 95\% CI
[$+0.005$, $+0.016$], distinguishable from zero;
Section~\ref{sec:uncertainty}) while eliminating the forcing-offset
assumption. The main remaining
absolute-temperature gap (modeled mean $\SI{40.77}{\degreeCelsius}$ vs.\
MODIS mean $\SI{47.56}{\degreeCelsius}$ on August 19) is attributable to
the model's simplified longwave radiative scheme rather than forcing, and
does not affect spatial pattern validation.

\paragraph{Diffusion coefficient and the baseline/intervention trade-off.}
\label{par:diffusion}
The single-stage calibration produces $D = 18.3\,\text{m}^2\,\text{s}^{-1}$,
which correctly reproduces the smooth spatial coherence of the baseline UHI
pattern but implies a canopy-cooling half-distance of 619\,m — substantially
longer than the 100--300\,m reported in field studies. The two-stage
calibration (Section~\ref{sec:cal_results}) reduces the half-distance to
363\,m by anchoring $D = 1.0\,\text{m}^2\,\text{s}^{-1}$ to the literature
and re-fitting the remaining coefficients, at a cost of 36\% in held-out MODIS
RMSE ($1.09 \to 1.48$\,°C) and significant narrowing of the physics advantage
over spatially-augmented regression ($r = +0.850 \to +0.747$). This quantified trade-off — baseline pattern
fidelity versus intervention spatial realism — is a fundamental constraint
of the single-parameter diffusion formulation: $D$ cannot simultaneously
reproduce the MODIS-observed spatial autocorrelation \emph{and} the
literature-observed local tree-cooling decay scale, because the two
observational targets imply different effective length scales.

The practical consequence for the targeted-zone ranking
(Section~\ref{sec:zonal}) is that rankings are likely directionally correct
but over-credit zones downwind of a planted patch; city-aggregate scenario
claims remain unaffected. Two model-structural changes would resolve the
conflict in future versions: (i) replacing the Fickian diffusion term with
an advection-diffusion formulation driven by ERA5 wind fields, which separates
spatial smoothing (advection) from physical heat conduction (diffusion), and
(ii) anisotropic diffusion coefficients distinguishing horizontal thermal
conduction within the substrate from atmospheric mixing above it.

\paragraph{Water-buffer parameterization.}
An earlier two-stage run using a simpler proportional-relaxation water
buffer produced an unphysical $\alpha_{\mathrm{wat}} = 0.49$ (literature:
0.06--0.10): with low $D$, water cells could not share cooling with
neighbors, and the optimizer raised absorption to compensate. This was
resolved by replacing the proportional-relaxation buffer with a proper
shallow mixed-layer model (explicit heat capacity ratio and latent-heat
flux term), which constrains $\alpha_{\mathrm{wat}}$ via its physical
bounds (0.04--0.18) and returned the optimized value to 0.17 —
within a factor of two of the literature range. The mixed-layer physics
is now the default; the remaining coefficient gap ($\alpha_{\mathrm{wat}}
= 0.17$ vs.\ literature 0.06--0.10) is modest and attributable to the
2021 WorldCover water mask not distinguishing shallow reservoirs from
open river channels.

\paragraph{Counterfactual scenarios.}
The counterfactual scenarios remain validated only against external literature
rather than against any direct observation of an implemented intervention in
Austin; the river greenway and densification scenarios have no literature
baseline at all. The canopy scenario applies a spatially uniform perturbation
and therefore cannot reproduce the threshold-dependent local cooling response
reported by \citet{ziter2019scale}. For the cool-roof scenario specifically,
the treatment-area-mean $\Delta T$ within the 3\,km downtown radius has now
been extracted directly from the scenario grid (Section~\ref{sec:lit_comparison}):
at the as-shipped albedo boost (0.35) it is \SI{-3.46}{\degreeCelsius}, above
the \SI{0.35}{}--\SI{1.2}{\degreeCelsius} literature band, but this reflects
scenario intensity rather than model error — the as-shipped boost is a more
aggressive retrofit than the literature studies measured. A re-run at a
literature-matched albedo boost (0.10) gives \SI{-0.97}{\degreeCelsius},
inside the band, and the two runs scale linearly in the albedo perturbation
to within rounding ($-0.97 \times 3.5 \approx -3.4$, vs.\ the as-shipped
\SI{-3.46}{\degreeCelsius}). We keep the as-shipped boost = 0.35 scenario as
the primary entry in Table~\ref{tab:scenarios} — preserving the existing
scenario configuration and prior \texttt{NUMBERS} entries rather than
reparameterizing the shipped default — and report both intensities side by
side in Table~\ref{tab:lit_comparison} so the literature check is made at
matched intensity. We consider this preferable to silently swapping the
shipped scenario to boost = 0.10: a reader auditing the as-shipped run
against Table~\ref{tab:lit_comparison} can see both numbers and the scaling
relationship between them, rather than a single re-parameterized figure
whose relationship to the originally reported \SI{-3.57}{}--\SI{-4.80}{\degreeCelsius}
peak-cell values (Table~\ref{tab:scenarios}) would otherwise be opaque.

\subsection{Evapotranspiration Formulation: Negative Result and Mechanism}
\label{sec:pm_negative}

The linear ET formulation under-predicts local canopy cooling by a factor
of $1.5$--$5\times$ relative to the empirical literature
(Section~\ref{sec:lit_comparison}), motivating an investigation of
Penman-Monteith (PM) ET as a physically grounded replacement. We implemented
canonical FAO-56 PM with LAI-weighted canopy conductance
($\mathrm{LAI} = 5 \cdot (1 - e^{-3 \cdot f_{\mathrm{veg}}})$),
calibrated under ERA5 forcing, and compared it against the legacy linear-ET
formulation across three diagnostic variants (Table~\ref{tab:pm_diagnostic}).

\begin{table}[ht]
  \centering
  \caption{Held-out test Pearson $r$ under four ET formulations, all using
  ERA5 forcing and two-stage diffusion anchor. None of the test days entered
  calibration.}
  \label{tab:pm_diagnostic}
  \begin{tabular}{lcc}
    \toprule
    ET formulation & Train $r$ & Test $r$ \\
    \midrule
    Legacy linear-ET + synthetic  & $+0.725$ & $+0.747$ \\
    Legacy linear-ET + ERA5       & $+0.714$ & $+0.735$ \\
    PM + ERA5 (LAI $K=3$, canonical) & $+0.496$ & $+0.418$ \\
    PM + ERA5 (LAI $K=1$, softer)    & $+0.567$ & $+0.515$ \\
    PM + ERA5 (linear LAI)           & $+0.548$ & $+0.493$ \\
    PM + ERA5 (LAI $= 10 \cdot f_{\mathrm{veg}}$) & $+0.484$ & $+0.411$ \\
    \bottomrule
  \end{tabular}
\end{table}

The diagnostic isolates two contributing mechanisms, roughly equal in size.
The first is \textbf{LAI saturation}: the exponential
$\mathrm{LAI} = 5(1 - e^{-3 f_{\mathrm{veg}}})$ reaches 84\% of its
maximum at $f_{\mathrm{veg}} = 0.5$, which collapses the ET gradient
across the top half of the vegetation-fraction range. Since Austin's
median vegetation fraction is 0.66 (Section~\ref{sec:data}), most city
cells sit in the flat tail where PM-ET is nearly invariant to land cover,
eliminating the spatial contrast that drives the UHI pattern. Softening to
linear LAI ($\mathrm{LAI} = k \cdot f_{\mathrm{veg}}$) recovers $\Delta r
\approx +0.09$ on the test set but does not close the gap.

The second, larger contributor is \textbf{structural VPD decoupling}:
in canonical FAO-56 PM, vapor pressure deficit depends only on
$(T_{\mathrm{air}}, T_{\mathrm{dew}})$ — atmospheric state variables
that are spatially uniform at ERA5's $\sim$30\,km resolution and identical
across every cell on a given timestep. This eliminates the surface-temperature
feedback that the linear ET term provides: in the linear formulation,
a hotter impervious cell drives larger $Q_{\mathrm{ET}}$ in neighboring
vegetated cells via $\max(T_i - T_{\mathrm{air}}, 0)$, amplifying spatial
contrast; in canonical PM, all cells see the same VPD regardless of local
surface temperature. Together, LAI saturation and VPD decoupling account
for the full $\Delta r \approx -0.33$ regression relative to the linear-ET
baseline, with no variant of the PM formulation recovering competitive
performance.

On the basis of this diagnostic, the ship configuration retains linear ET.
The negative PM result is not a failure of the Penman-Monteith framework
per se: PM is designed to estimate actual evapotranspiration from a
well-watered reference surface, not to drive spatial surface-temperature
contrasts in a simplified 2D slab model. The correct path to a more
physical ET formulation in this architecture is a \emph{big-leaf bulk PM}
with a $T_{\mathrm{surf}}$-dependent VPD term (replacing
$T_{\mathrm{air}}$ with $T_i$ in the saturation vapor pressure
calculation), which preserves the surface-temperature feedback while
removing the linear-ET simplification. This is identified as the next
ET model-structural change but is outside the scope of this paper.

\subsection{Comparison to Related Approaches}

This model occupies a distinct position relative to both WRF-type models
and ML-based approaches. Against WRF: we sacrifice spatial resolution,
atmospheric coupling, and absolute forcing accuracy, but gain zero setup
cost, city-scale runs in seconds, and a counterfactual engine. Against
ML models: we sacrifice the ability to learn from data at scale, but gain
causal physical interpretability and the ability to evaluate interventions
outside the observed distribution. This positioning suggests the model
is most suited for: rapid exploratory scenario analysis in early-stage
planning; educational contexts where a fast, transparent model is
preferable to a black box; and as a computationally inexpensive
``first-pass'' filter before running a small number of high-fidelity
WRF simulations on the most promising intervention designs.

% ── 7. Conclusion ────────────────────────────────────────────────────────────
\section{Conclusion}
\label{sec:conclusion}

We have presented \textsc{austin-twin}, a lightweight, fully reproducible
digital twin for urban heat island simulation and counterfactual planning.
The ship configuration — two-stage diffusion-anchored calibration, mixed-layer
water physics, and per-date ERA5 atmospheric forcing — achieves $r = +0.758$
and RMSE = \SI{1.42}{\degreeCelsius} on held-out MODIS Aqua land surface
temperature observations using only open, zero-authentication data sources,
with every calibrated coefficient in its published physical range. A
multi-city transfer experiment on Phoenix and Miami identifies WorldCover
land-cover separability — not climate distance — as the primary determinant
of portability, and introduces a data-ceiling control (Pearson correlation
between MODIS LST and impervious fraction) that separates genuine model
failure from sensor-driven class-separability limits; this diagnostic reverses
the apparent interpretation of the Phoenix result and is, as argued in
Section~\ref{sec:novelty}, the paper's most transferable methodological
contribution. A documented negative result for canonical Penman-Monteith ET
isolates the two causal mechanisms (LAI saturation and VPD decoupling) and
defines the correct architectural path forward. The framework is available
publicly at [repository URL] and requires only a standard Python environment
to run.

Future work will pursue two open threads identified by the diagnostic
experiments in this paper. First, a big-leaf bulk PM ET formulation with
$T_{\mathrm{surf}}$-dependent VPD preserves the surface-temperature
feedback eliminated by canonical PM while removing the linear-ET
simplification; this is the natural next ET model-structural step
(Section~\ref{sec:pm_negative}). Second, iterating the Stage-1/Stage-2
diffusion anchor to convergence, or extending the $D$-sweep below
$1.0\,\text{m}^2\,\text{s}^{-1}$, would close the residual 70--160\,m
gap between modeled and literature-reported canopy-cooling half-distances.
Beyond these: re-calibration at 100\,m resolution to close the
highway-corridor bias; extension to cities where WorldCover separability
is high (Houston, Seattle) as further zero-shot transfer tests; investigation
of whether augmenting WorldCover with an irrigation-layer channel would
raise Phoenix's data ceiling; and re-calibration at 100\,m resolution
to close the highway-corridor bias.

% ── Acknowledgments ──────────────────────────────────────────────────────────
\section*{Acknowledgments}
The author thanks [mentor name] for guidance throughout this project.
MODIS data were obtained from the Microsoft Planetary Computer STAC API.
ESA WorldCover data were produced by the ESA WorldCover project with funding
by ESA. OpenStreetMap data are copyright OpenStreetMap contributors and
available under the Open Database License. ERA5 data were generated using
Copernicus Climate Change Service information; neither the European Commission
nor ECMWF is responsible for any use that may be made of this information.

% ── Data Availability ─────────────────────────────────────────────────────────
\section*{Data and Code Availability}
All code, calibration outputs, and figure-generation scripts are available at
\url{https://github.com/[username]/austin-twin}. All data sources are freely
accessible without registration as described in Section~\ref{sec:data}.
Generated using Copernicus Climate Change Service information [2024].

% ── References ───────────────────────────────────────────────────────────────
% Note: this bibliography is hand-written (no \bibliography{} / BibTeX run).
% \bibliographystyle has no effect on formatting here (nothing invokes
% BibTeX), but is kept for documentation; natbib's author-year mode is set
% directly via the [authoryear] package option above, and each \bibitem
% uses the comma-separated "[Author, Year]" bracket form natbib's
% hand-bibliography parser expects (see \NAT@apalk in natbib.sty) — a
% "[Author(Year)]" bracket with no comma silently forces numeric citations
% at \end{document}, which was the cause of the natbib "not compatible
% with author-year citations" error.
\bibliographystyle{plainnat}
\begin{thebibliography}{99}

\bibitem[Bhamjee et al., 2024]{bhamjee2024detection}
Bhamjee, M., et al.\ (2024).
Detection and characterization of urban heat islands with machine learning.
In \textit{IGARSS 2024 -- IEEE International Geoscience and Remote Sensing
Symposium}, pp.\ 1693--1699.

\bibitem[Bueno et al., 2013]{bueno2013urban}
Bueno, B., Norford, L., Hidalgo, J., and Pigeon, G.\ (2013).
The urban weather generator.
\textit{Journal of Building Performance Simulation}, 6(4), 269--281.

\bibitem[Hersbach et al., 2020]{hersbach2020era5}
Hersbach, H., et al.\ (2020).
The ERA5 global reanalysis.
\textit{Quarterly Journal of the Royal Meteorological Society},
146(730), 1999--2049.
\url{https://doi.org/10.1002/qj.3803}.

\bibitem[HeatMat, 2025]{heatmat2025}
[Authors].\ (2025).
HeatMat: Simulation of city material impact on urban heat island effect.
\textit{arXiv preprint arXiv:2601.22796}.

\bibitem[Hoang and Nguyen, 2025]{hoang2025geospatial}
Hoang, N.D., and Nguyen, Q.L.\ (2025).
Geospatial analysis and machine learning framework for urban heat island
intensity prediction: Natural gradient boosting and deep neural network
regressors with multisource remote sensing data.
\textit{Sustainability}, 17, 4287.

\bibitem[Krayenhoff et al., 2021]{krayenhoff2021albedo}
Krayenhoff, E.S., et al.\ (2021).
Cooling hot cities: a systematic and critical review of the numerical
modelling literature.
\textit{Environmental Research Letters}, 16(5), 053007.

\bibitem[Kusaka et al., 2001]{kusaka2001simple}
Kusaka, H., Kondo, H., Kikegawa, Y., and Kimura, F.\ (2001).
A simple single-layer urban canopy model for atmospheric models:
Comparison with multi-layer and slab models.
\textit{Boundary-Layer Meteorology}, 101(3), 329--358.

\bibitem[Lyu et al., 2022]{lyu2022integrated}
Lyu, F., Wang, S., Han, S.Y., Catlett, C., and Wang, S.\ (2022).
An integrated CyberGIS and machine learning framework for fine-scale
prediction of urban heat island using satellite remote sensing and urban
sensor network data.
\textit{Urban Informatics}, 1, 6.

\bibitem[Masson, 2000]{masson2000physically}
Masson, V.\ (2000).
A physically-based scheme for the urban energy budget in atmospheric models.
\textit{Boundary-Layer Meteorology}, 94(3), 357--397.

\bibitem[Microsoft, 2023]{microsoft_pc}
Microsoft.\ (2023).
Planetary Computer STAC API.
\url{https://planetarycomputer.microsoft.com}.

\bibitem[Mitraka et al., 2025]{mitraka2025prototypes}
Mitraka, Z., Lantzanakis, G., Panagiotakis, E., et al.\ (2025).
Prototypes for enhanced urban heat monitoring leveraging Copernicus
services and satellite data.
\textit{Discover Cities}, 2, 14.

\bibitem[Oke, 1982]{oke1982energetic}
Oke, T.R.\ (1982).
The energetic basis of the urban heat island.
\textit{Quarterly Journal of the Royal Meteorological Society},
108(455), 1--24.

\bibitem[Santamouris, 2014]{santamouris2014cool}
Santamouris, M.\ (2014).
Cooling the cities -- a review of reflective and green roof mitigation
technologies to fight heat island and improve comfort in urban environments.
\textit{Solar Energy}, 103, 682--703.

\bibitem[Storn and Price, 1997]{storn1997differential}
Storn, R., and Price, K.\ (1997).
Differential evolution -- a simple and efficient heuristic for global
optimization over continuous spaces.
\textit{Journal of Global Optimization}, 11(4), 341--359.

\bibitem[Voogt and Oke, 2003]{voogt2003thermal}
Voogt, J.A., and Oke, T.R.\ (2003).
Thermal remote sensing of urban climates.
\textit{Remote Sensing of Environment}, 86(3), 370--384.

\bibitem[Wang and Li, 2022]{wang2022berlin}
Wang, X., and Li, H.\ (2022).
The effectiveness of cool and green roofs in mitigating urban heat island
and improving human thermal comfort.
\textit{Building and Environment}, 217, 109082.

\bibitem[World Resources Institute, 2026]{wri2026cooling}
Wesley, E.J., Mackres, E., and Pool, J.\ (2026).
Cooling potential of urban trees.
\textit{World Resources Institute Insights}.
\url{https://www.wri.org/insights/urban-trees-cooling-potential}.

\bibitem[Zanaga et al., 2022]{zanaga2022esa}
Zanaga, D., et al.\ (2022).
ESA WorldCover 10 m 2021 v200.
\textit{Zenodo}. \url{https://doi.org/10.5281/zenodo.7254221}.

\bibitem[Zhou et al., 2025]{zhou2025tree}
Zhou, W., et al.\ (2025).
Increasing tree canopy lowers urban air temperature by up to 1.5 °C in
heat-prone areas.
\textit{npj Urban Sustainability}, 5, 92.

\bibitem[Ziter et al., 2019]{ziter2019scale}
Ziter, C.D., Pedersen, E.J., Kucharik, C.J., and Turner, M.G.\ (2019).
Scale-dependent interactions between tree canopy cover and impervious
surfaces reduce daytime urban heat during summer.
\textit{Proceedings of the National Academy of Sciences}, 116(15), 7575--7580.

\end{thebibliography}

\end{document}
